\documentclass{situationist-preprint}

\usepackage{longtable,array}
\newenvironment{errlist}{\small\setlength{\LTpre}{2pt}\setlength{\LTpost}{4pt}%
\begin{longtable}{@{}>{\raggedleft}p{1.2cm}@{\hspace{0.8em}}>{\raggedright}p{3.6cm}@{\hspace{0.8em}}>{\raggedright\arraybackslash}p{\dimexpr\linewidth-4.8cm-1.6em\relax}@{}}}%
{\end{longtable}}
\newenvironment{skeltab}{\small\setlength{\LTpre}{2pt}\setlength{\LTpost}{8pt}%
\begin{longtable}{@{}r@{\hspace{0.8em}}p{\dimexpr\linewidth-3.6cm\relax}@{\hspace{0.8em}}r@{}}}%
{\end{longtable}}

\title{Errata to \enquote{On the Good}, Volume IV}
\msid{ETH-01-e}
\msfamily{ETH-01}
\msdiscipline{Normative and Applied Ethics}
\mstype{Erratum}
\mskeywords{errata; fairness; claims; lotteries; priority; isolation test}

\begin{document}
\maketitle

\begin{abstract}
We list the corrections to Volume~IV of \emph{On the Good} (ETH-01), \emph{Fairness}, pages 1,161 to 1,559. Six are substantive: a restatement of the isolation pair for fairness that omits one of the respects in which the two outcomes agree, a proposition on weighted lotteries stated without its restriction to indivisible goods, two citations on which arguments rely, a misnumbered reference to Volume~I in the proof that fairness is not a component of well-being, and a slip in the statement of what Volume~VI establishes. For each we give the printed and the corrected text and show that the results that use it stand. The remaining corrections, about nine hundred and fifty, are misprints and slips of reference, listed by chapter and page. None affects the conclusion of the volume, that fairness passes the isolation test and is an irreducible component of the good.
\end{abstract}

\section{Preface}

Volume~IV of \emph{On the Good} applies the method of Volume~I (ETH-01-a) to fairness, the fifth of the seven components of the good. It argues that fairness passes the isolation test of Definition~4.1 of Volume~I, that it cannot be reduced to equality, to the priority of the worse off or to desert, and that it is a good of outcomes which is not a component of anyone's well-being. Pagination is continuous across the six volumes, and page numbers below are those of Volume~IV. Chapters are numbered afresh in the volume. Definitions, lemmas, propositions, theorems, corollaries and remarks share one counter within each chapter, as in Volume~I, and displayed formulas are numbered (4.1), (4.2), and so on within each chapter. The chapters and their page ranges are given in Section~\ref{s:skeleton}.

The list was compiled in three ways: by checking every cross-reference against the item it names, and every page reference against the page of that item; by checking every citation against the bibliography of the volume and every bibliographic entry against the published work; and from readers' reports and our own rereading, which account for most of the typographical corrections. The first two checks were complete. We do not claim completeness for typographical errors in running text.

We distinguish two kinds of correction. A correction is \emph{substantive} if it changes what a numbered statement or a displayed condition says, or which work an argument relies on. The six substantive corrections are given in full in Section~\ref{s:substantive}, as E1 to E6, each with the results that use the corrected item; at its place in the list by chapter a pointer refers to it. All other corrections, in Section~\ref{s:list}, are misprints of words, symbols and numbers, omitted or repeated words, wrong cross-references and page references, and errors in citations, in the bibliography and in the indices. None of them changes what a numbered result states. Entries inside the statement or proof of a numbered result are marked with a dagger (\dag); for each we have checked that the proof is valid with the correction made and that the corrected reading is the one the proof uses. Where the correct form is printed elsewhere in the chapter, the entry says where.

No correction affects the conclusions of the volume or of the work. Theorem~2.6, that fairness passes the isolation test, stands as printed (E1), and with it the contribution of Volume~IV to Assumption~5.1 of Volume~I, the conjunction of the seven isolation results from which Theorem~5.2 of Volume~I derives the plurality of the good. Theorem~10.4, that fairness is not a component of anyone's well-being, stands (E5). The volume makes no claim about the completeness of the list beyond reporting, correctly once E6 is made, that Volume~VI identifies completeness with the supervenience of value on all seven components together.

Entries have the form \enquote{\textit{for} X \textit{read} Y}. Lines are counted from the top of the text block, without the running head; \enquote{from foot} counts upward from the last line of text; \enquote{n.} is a footnote. Entries in the bibliography are located by the heading of the entry, and entries in the indices by the heading under which the page is listed. The list contains the following numbers of entries.

\begin{center}\small
% generated by generate.py -- do not edit
\begin{tabular}{@{}lr@{}}
\toprule
Typographical errors (spelling, running heads) & 615 \\
Wrong word or phrase & 22 \\
Omitted or repeated words & 28 \\
Misprinted formulas and symbols & 114 \\
Cross-references and page references & 197 \\
Citations in the text & 95 \\
Bibliography & 28 \\
Indices & 49 \\
\midrule
Total & 1148 \\
\bottomrule
\end{tabular}

\end{center}

Three entries show the kinds of error that the cross-checks found. On p.~1,163 the Contents give the first page of Chapter~6 as 1,222 for 1,322. In Chapter~2, p.~1,194, display (2.3), which says that the chances of a lottery sum to one, sums over the wrong index, so that as printed it states that a single candidate's chance, counted once for each candidate, is one; the correct form is printed in (2.5) on p.~1,198. In the Index of Names, p.~1,550, the reference to Crisp is printed as 1,434 for 1,344, a page of Chapter~6, where the sufficiency view is discussed.

\section{Chapters of Volume IV}\label{s:skeleton}

% generated by generate.py -- do not edit
\begin{skeltab}
 & Front matter & 1,161--1,164 \\
1 & Fairness as a component of the good & 1,165--1,185 \\
2 & The isolation test for fairness & 1,186--1,218 \\
3 & Fairness and equality & 1,219--1,247 \\
4 & Lotteries and indivisible goods & 1,248--1,283 \\
5 & Claims and their satisfaction & 1,284--1,321 \\
6 & Priority, sufficiency and levelling down & 1,322--1,360 \\
7 & Procedural and outcome fairness & 1,361--1,392 \\
8 & Fairness across generations & 1,393--1,420 \\
9 & Fairness and desert & 1,421--1,446 \\
10 & Fairness is not a component of well-being & 1,447--1,477 \\
11 & Aggregation and the numbers & 1,478--1,507 \\
12 & Summary: fairness in the list of goods & 1,508--1,520 \\
 & Bibliography & 1,521--1,548 \\
 & Index of names & 1,549--1,552 \\
 & Index of subjects & 1,553--1,559 \\
\end{skeltab}


\section{Substantive errata}\label{s:substantive}

\paragraph{E1. Chapter~2, display (2.4), p.~1,197.} Display (2.4) states the isolation pair for fairness: two outcomes with the same total well-being, in one of which the distribution results from a lottery to which everyone consented, and in the other from a procedure that favored some arbitrarily. \textit{After} \enquote{the same total well-being} \textit{insert} \enquote{and the same distribution of well-being, knowledge, achievement, relationships, beauty and autonomy}. As printed, the two outcomes may differ in the distribution of well-being, so that a preference for the first could be a preference for equality of outcome and not for fairness. The pair is stated correctly in Volume~I, Chapter~4 (ETH-01-a, Section~4), and in the paragraph that follows the display on p.~1,198, which says that the two outcomes are alike in every component other than fairness. Theorem~2.6 (p.~1,203), that fairness passes the isolation test, is proved from the pair as stated on p.~1,198, and it stands.

\paragraph{E2. Chapter~4, Proposition~4.5, p.~1,268.} The proposition is printed: \enquote{A lottery whose chances are proportional to the strengths of the candidates' claims satisfies their claims fairly.} \textit{After} \enquote{A lottery} \textit{insert} \enquote{for an indivisible good}. When the good can be divided, fairness requires that the claims be satisfied in proportion to their strength by dividing the good, and a lottery, which gives the whole good to one candidate, satisfies the claims of the others not at all \citep{Broome1990}. The proof (pp.~1,268--1,271) assumes indivisibility in its first line, and the chapter title restricts the subject accordingly. Proposition~4.5 is used in Theorem~4.8 (p.~1,276), on the allocation of a single organ among candidates with claims of different strength, where the good is indivisible. Theorem~4.8 stands, and so do the uses of the proposition in Chapter~11.

\paragraph{E3. Chapter~4, p.~1,252, l.~9, and Bibliography, p.~1,528, entry Diamond (1967).} Diamond's objection to Harsanyi's aggregation theorem, that a social ranking indifferent between a lottery that gives the good to one person for certain and a lottery that gives each of two persons an equal chance ignores the fairness of the procedure, is attributed to a paper in \emph{Econometrica}. \textit{For} Econometrica \textit{read} Journal of Political Economy, 75(5), 765--766, in both places \citep{Diamond1967,Harsanyi1955}. The year and the title of the entry are correct. The discussion in Chapter~4 reproduces the objection correctly and is unaffected.

\paragraph{E4. Chapter~6, p.~1,349, l.~3 from foot.} The levelling-down objection is attributed to \enquote{Temkin (1993)}. \textit{For} the levelling-down objection of Temkin (1993) \textit{read} the levelling-down objection of Parfit (1997). The objection, that on the telic egalitarian view an outcome in which the better off are brought down to the level of the worse off is in one respect better although it is better for no one, was named and pressed by \citet{Parfit1997}; \citet{Temkin1993} is the egalitarian reply to it, cited correctly on p.~1,351. The chapter's own reply to the objection, which distinguishes the value of fairness from the value of equality, is unaffected.

\paragraph{E5. Chapter~10, proof of Theorem~10.4, p.~1,459.} Theorem~10.4 states that fairness is not a component of anyone's well-being: there are two outcomes alike in the well-being of every individual that differ in fairness and in value. The proof cites \enquote{Corollary~5.3 of Volume~I}. \textit{For} Corollary~5.3 \textit{read} Corollary~5.4. Corollary~5.3 of Volume~I states that the value of an outcome is not a function of any single component; the step needs Corollary~5.4, that value is not a function of well-being alone, which, as Volume~I notes, holds also when well-being is represented by the vector of individual welfare levels. That is the form the proof applies. Theorem~10.4 stands as printed.

\paragraph{E6. Chapter~12, p.~1,515, l.~18.} The summary states that Volume~VI \enquote{shows that the value of outcomes supervenes on the six components}. \textit{For} six \textit{read} seven. Completeness of the list, the subject of Volume~VI, is the supervenience of value on all seven components together, as ETH-01-a states; supervenience on six would contradict Theorem~5.2 of Volume~I, by which value does not supervene on any set that omits one of the seven. Nothing in Volume~IV depends on the sentence.

\section{Errata by chapter}\label{s:list}
% generated by generate.py -- do not edit
\subsection*{Front matter \hfill \normalfont\small pp.~1,161--1,164}
\begin{errlist}
1,163 & Contents, Chapter~6 & \textit{for} 1,222 \textit{read} 1,322 \\
\end{errlist}

\subsection*{Chapter~1. Fairness as a component of the good \hfill \normalfont\small pp.~1,165--1,185}
\begin{errlist}
1,166 & l.~16 & \textit{for} furhter \textit{read} further \\
1,166 & l.~30 & \textit{for} Definition~4.1 of Volume~I \textit{read} Definition~3.1 of Volume~I \\
1,166 & n.~14 & \textit{for} Rawls (1972) \textit{read} Rawls (1971) \\
1,167 & l.~37 & \textit{delete} the second occurrence of \emph{is} \\
1,168 & l.~14 & \textit{for} contrbution \textit{read} contribution \\
1,168 & l.~18 & \textit{for} excepsion \textit{read} exception \\
1,169 & l.~19 & \textit{for} absrtact \textit{read} abstract \\
1,169 & l.~28 & \textit{for} Hookr (2005) \textit{read} Hooker (2005) \\
1,169 & l.~37 & \textit{for} elswehere \textit{read} elsewhere \\
1,170 & l.~5 & \textit{for} dleiberate \textit{read} deliberate \\
1,170 & l.~16 & \textit{for} Hooker (2006) \textit{read} Hooker (2005) \\
1,170 & l.~17 & \textit{for} elsewhree \textit{read} elsewhere \\
1,170 & l.~25 & \textit{delete} the second occurrence of \emph{a} \\
1,170 & l.~27 & \textit{after} lottery \textit{insert} with chances proportional to the claims \\
1,170 & l.~31 & \textit{delete} the second occurrence of \emph{that} \\
1,170 & l.~38 & \textit{for} mnaifest \textit{read} manifest \\
1,171 & l.~4 & \textit{for} sufficiant \textit{read} sufficient \\
1,171 & l.~17 & \textit{for} effectivley \textit{read} effectively \\
1,171 & l.~21 & \textit{for} Moore (1903) \textit{read} Broome (1990) (the claims account of fairness) \\
1,171 & l.~22 & \textit{for} staisfaction \textit{read} satisfaction \\
1,171 & l.~26 & \textit{for} Volume~I \textit{read} Volume~II (well-being) \\
1,172 & l.~6 & \textit{after} claims \textit{insert} to the good \\
1,172 & l.~9 & \textit{for} an equal-chance lottery is unfair when the claims are equal \textit{read} an equal-chance lottery is fair when the claims are equal \\
1,172 & l.~12 & \textit{for} soltion \textit{read} solution \\
1,172 & l.~21 & \textit{for} Broome (1989) \textit{read} Broome (1990) \\
1,173 & l.~7 & \textit{for} epsecially \textit{read} especially \\
1,173 & l.~13 & \textit{for} nohting \textit{read} nothing \\
1,173 & l.~19 & \textit{for} nosion \textit{read} notion \\
1,173 & l.~26 & \textit{for} Volume~I, Chapter~1 \textit{read} Volume~I, Chapter~3 (idealized reflective endorsement) \\
1,173 & l.~30 & \textit{for} evdent \textit{read} evident \\
1,173 & l.~40 & \textit{for} reefrence \textit{read} reference \\
1,174 & l.~2 & \textit{for} Volume~IV, Chapter~10 \textit{read} Volume~IV, Chapter~11 (the numbers) \\
1,174 & l.~4 from foot & \textit{for} adidtional \textit{read} additional \\
1,175 & l.~4 & \textit{for} neihter \textit{read} neither \\
1,175 & l.~7 & \textit{for} objectve \textit{read} objective \\
\dag\,1,175 & proof of Lemma~1.2, l.~9 & \textit{for} Definition~3.1 of Volume~I \textit{read} Definition~4.1 of Volume~I \\
1,175 & l.~18 & \textit{for} generrally \textit{read} generally \\
1,175 & l.~20 & \textit{for} outcome fairness \textit{read} procedural fairness \\
1,175 & l.~23 & \textit{for} ambguity \textit{read} ambiguity \\
1,175 & l.~33 & \textit{for} p.~1,168 \textit{read} p.~1,167 (Definition~1.1) \\
1,176 & l.~7 & \textit{for} exsitence \textit{read} existence \\
1,176 & l.~9 & \textit{for} consderation \textit{read} consideration \\
1,176 & l.~22 & \textit{for} aruged \textit{read} argued \\
1,176 & l.~32 & \textit{delete} the second occurrence of \emph{is} \\
1,177 & l.~38 & \textit{for} langugae \textit{read} language \\
1,178 & l.~1 from foot & \textit{for} qeustion \textit{read} question \\
1,179 & l.~4 & \textit{for} suffciient \textit{read} sufficient \\
1,179 & l.~5 & \textit{for} Broomme (1990) \textit{read} Broome (1990) \\
1,180 & l.~5 & \textit{for} caapcity \textit{read} capacity \\
1,180 & l.~6 & \textit{for} argeud \textit{read} argued \\
1,180 & l.~25 & \textit{for} rceeive \textit{read} receive \\
1,181 & l.~5 & \textit{for} motiavtion \textit{read} motivation \\
1,181 & l.~30 & \textit{for} Volume~I, Chapter~8 \textit{read} Volume~I, Chapter~6 (incommensurability and parity) \\
1,181 & l.~32 & \textit{for} Sigdwick (1907) \textit{read} Sidgwick (1907) \\
1,182 & l.~5 & \textit{for} Volume~I, Chapter~4 \textit{read} Volume~I, Chapter~6 (incommensurability and parity) \\
1,183 & l.~3 & \textit{for} scucessful \textit{read} successful \\
1,183 & l.~20 & \textit{for} p.~1,181 \textit{read} p.~1,180 (Definition~1.3) \\
1,184 & l.~30 & \textit{for} presumbly \textit{read} presumably \\
1,184 & n.~15 & \textit{for} Moore (1902) \textit{read} Moore (1903) \\
1,185 & l.~8 & \textit{for} equivlent \textit{read} equivalent \\
1,185 & l.~22 & \textit{for} succed \textit{read} succeed \\
1,185 & l.~27 & \textit{for} reqiurement \textit{read} requirement \\
1,185 & l.~6 from foot & \textit{for} respectievly \textit{read} respectively \\
\end{errlist}

\subsection*{Chapter~2. The isolation test for fairness \hfill \normalfont\small pp.~1,186--1,218}
\begin{errlist}
1,187 & l.~17 & \textit{for} $y_{-F} = x_{-F} \vee y_{F} \neq x_{F} \wedge V(y) \neq V(x)$ \textit{read} $y_{-F} = x_{-F} \wedge y_{F} \neq x_{F} \wedge V(y) \neq V(x)$ \\
1,188 & l.~15 & \textit{for} correspnoding \textit{read} corresponding \\
1,188 & l.~25 & \textit{for} Moore (1904) \textit{read} Moore (1903) \\
1,188 & display (2.1) & \textit{for} $F(z) < F(x) \wedge z_{-F} = x_{-F} \to V(z) > V(x)$ \textit{read} $F(z) > F(x) \wedge z_{-F} = x_{-F} \to V(z) > V(x)$ \\
1,189 & l.~17 & \textit{for} differnt \textit{read} different \\
1,190 & l.~4 & \textit{for} occurence \textit{read} occurrence \\
1,190 & l.~17 & \textit{for} Theorem~5.2 of Volume~I (the isolation results) \textit{read} Assumption~5.1 of Volume~I (the isolation results) \\
1,190 & l.~30 & \textit{for} the stronger claim is satisfied only in part \textit{read} the weaker claim is satisfied only in part \\
1,191 & l.~5 & \textit{for} tendncy \textit{read} tendency \\
1,191 & l.~7 & \textit{for} Volume~I, Chapter~5 \textit{read} Volume~I, Chapter~3 (idealized reflective endorsement) \\
1,191 & l.~12 & \textit{for} silghtly \textit{read} slightly \\
1,191 & l.~36 & \textit{for} prolbem \textit{read} problem \\
1,191 & l.~4 from foot & \textit{for} irreducable \textit{read} irreducible \\
1,192 & l.~1 & \textit{for} Moroe (1903) \textit{read} Moore (1903) \\
1,192 & l.~4 & \textit{for} diferent \textit{read} different \\
1,192 & l.~9 & \textit{for} doctrnie \textit{read} doctrine \\
1,192 & l.~13 & \textit{for} inded \textit{read} indeed \\
1,192 & l.~15 & \textit{for} consideralbe \textit{read} considerable \\
1,192 & l.~20 & \textit{for} modset \textit{read} modest \\
1,192 & display (2.2) & \textit{for} $c_{j} > c_{j} \to \pi_{i} > \pi_{j}$ \textit{read} $c_{i} > c_{j} \to \pi_{i} > \pi_{j}$ \\
1,192 & l.~31 & \textit{for} $\sum_{j} s_{k} \leq S$ \textit{read} $\sum_{k} s_{k} \leq S$ \\
1,192 & l.~36 & \textit{for} ultimattely \textit{read} ultimately \\
1,193 & l.~1 & \textit{for} restricsion \textit{read} restriction \\
1,193 & l.~1 from foot & \textit{for} Ross (1931) \textit{read} Ross (1930) \\
1,194 & l.~4 & \textit{for} Definition~4.1 of Volume~I (the good) \textit{read} Definition~3.2 of Volume~I (the good) \\
1,194 & l.~8 & \textit{for} dessirable \textit{read} desirable \\
1,194 & display (2.3) & \textit{for} $\sum_{j} \pi_{i} = 1$ \textit{read} $\sum_{i} \pi_{i} = 1$; correctly printed in (2.5), p.~1,198 \\
1,195 & l.~1 & \textit{for} advnatage \textit{read} advantage \\
1,195 & l.~2 & \textit{for} virtualy \textit{read} virtually \\
1,195 & l.~12 & \textit{for} distintcion \textit{read} distinction \\
1,195 & l.~5 from foot & \textit{for} judgmeent \textit{read} judgement \\
1,195 & l.~40 & \textit{for} rgearded \textit{read} regarded \\
1,196 & l.~9 & \textit{for} propotrion \textit{read} proportion \\
1,196 & l.~14 & \textit{for} $\sum_{i} \pi_{j} = 1$ \textit{read} $\sum_{i} \pi_{i} = 1$ \\
1,196 & l.~17 & \textit{for} Assumption~5.1 of Volume~I (plurality) \textit{read} Theorem~5.2 of Volume~I (plurality) \\
1,197 & display (2.4) & \textit{see} Erratum~E1 \\
1,199 & l.~11 & \textit{for} Volume~I, Chapter~1 \textit{read} Volume~I, Chapter~2 (the concept of the good) \\
1,199 & display (2.7) & \textit{for} $\pi_{j} = c_{i} \Big/ \sum_{i} c_{i}$ \textit{read} $\pi_{j} = c_{j} \Big/ \sum_{i} c_{i}$ \\
1,199 & display (2.6) & \textit{for} $F(y) > F(x) \wedge y_{-F} = x_{-F} \leftrightarrow V(y) > V(x)$ \textit{read} $F(y) > F(x) \wedge y_{-F} = x_{-F} \to V(y) > V(x)$; the formula is used in its correct form on p.~1,203 \\
1,200 & l.~4 & \textit{for} combinasion \textit{read} combination \\
1,200 & l.~7 & \textit{for} commtited \textit{read} committed \\
1,200 & l.~13 & \textit{for} sequance \textit{read} sequence \\
1,200 & l.~14 & \textit{for} alhtough \textit{read} although \\
\dag\,1,200 & proof of Corollary~2.4, l.~21 & \textit{for} Definition~2.2 \textit{read} Definition~2.1 \\
\dag\,1,200 & proof of Corollary~2.4, (2.8) & \textit{for} $z_{-F} = x_{-F} \vee z_{F} \neq x_{F} \wedge V(z) \neq V(x)$ \textit{read} $z_{-F} = x_{-F} \wedge z_{F} \neq x_{F} \wedge V(z) \neq V(x)$; correctly printed in (2.10), p.~1,203 \\
\dag\,1,200 & proof of Corollary~2.4, l.~30 & \textit{for} Definition~2.3 \textit{read} Definition~2.2 \\
1,201 & l.~2 & \textit{for} combinaiton \textit{read} combination \\
\dag\,1,201 & proof of Lemma~2.5, (2.9) & \textit{for} $\pi_{j} = c_{k} \Big/ \sum_{k} c_{k}$ \textit{read} $\pi_{j} = c_{j} \Big/ \sum_{k} c_{k}$ \\
1,201 & l.~12 & \textit{for} supervenes on every set of components that omits fairness \textit{read} does not supervene on any set of components that omits fairness \\
1,201 & l.~28 & \textit{for} complte \textit{read} complete \\
1,202 & l.~11 & \textit{for} Definition~3.2 of Volume~I (the isolation test) \textit{read} Definition~4.1 of Volume~I (the isolation test) \\
1,202 & l.~31 & \textit{for} Volume~I, Chapter~3 \textit{read} Volume~I, Chapter~5 (the formal argument for plurality) \\
1,204 & l.~13 & \textit{for} Volume~II \textit{read} Volume~III (relationships) \\
1,205 & l.~7 & \textit{for} primmarily \textit{read} primarily \\
1,205 & l.~17 & \textit{for} a claim is a reason to give the good to whoever can use it best \textit{read} a claim is a reason owed to the candidate herself \\
\dag\,1,206 & proof of Proposition~2.7, (2.11) & \textit{for} $c_{i} < c_{j} \to \pi_{i} > \pi_{j}$ \textit{read} $c_{i} > c_{j} \to \pi_{i} > \pi_{j}$ \\
1,207 & l.~10 & \textit{for} Volume~IV, Chapter~10 \textit{read} Volume~IV, Chapter~8 (fairness across generations) \\
\dag\,1,207 & proof of Proposition~2.7, (2.12) & \textit{for} $x \succ z \iff V(x) < V(z)$ \textit{read} $x \succ z \iff V(x) > V(z)$; the formula is used in its correct form on p.~1,208 \\
1,207 & l.~6 from foot & \textit{for} prevoius \textit{read} previous \\
1,208 & l.~3 & \textit{for} unviersal \textit{read} universal \\
1,208 & l.~6 & \textit{for} Volume~IV, Chapter~2 \textit{read} Volume~IV, Chapter~3 (fairness and equality) \\
1,208 & l.~9 & \textit{for} Volume~I \textit{read} Volume~III (achievement) \\
1,208 & l.~15 & \textit{for} improtance \textit{read} importance \\
1,208 & l.~18 & \textit{for} p.~1,188 \textit{read} p.~1,187 (Definition~2.1) \\
1,208 & l.~20 & \textit{for} Volume~IV, Chapter~9 \textit{read} Volume~IV, Chapter~7 (procedural fairness) \\
\dag\,1,208 & proof of Theorem~2.10, (2.13) & \textit{for} $W_{k}(z) = W_{i}(x) \text{ for all } k \wedge F(z) \neq F(x)$ \textit{read} $W_{k}(z) = W_{k}(x) \text{ for all } k \wedge F(z) \neq F(x)$; the formula is used in its correct form on p.~1,210 \\
1,208 & n.~13 & \textit{for} Ross (1929) \textit{read} Ross (1930) \\
\dag\,1,209 & proof of Theorem~2.10, l.~11 & \textit{for} Definition~2.11 \textit{read} Definition~2.9; the step uses Definition~2.9 (p.~1,207) \\
1,209 & l.~34 & \textit{delete} the second occurrence of \emph{that} \\
1,209 & l.~40 & \textit{for} advantge \textit{read} advantage \\
1,210 & l.~14 & \textit{for} Ross (1930) \textit{read} Broome (1990) (the claims account of fairness) \\
1,210 & l.~19 & \textit{for} convicsion \textit{read} conviction \\
1,210 & l.~7 from foot & \textit{for} threeby \textit{read} thereby \\
1,210 & l.~35 & \textit{for} prleiminary \textit{read} preliminary \\
1,211 & l.~5 & \textit{for} Volume~IV \textit{read} Volume~III (achievement) \\
1,211 & l.~30 & \textit{for} p.~1,308 \textit{read} p.~1,208 (Theorem~2.10) \\
1,211 & l.~32 & \textit{for} $\forall y\,\forall z\,(y|_{S} = z|_{S} \leftrightarrow V(y) = V(z))$ \textit{read} $\forall y\,\forall z\,(y|_{S} = z|_{S} \to V(y) = V(z))$ \\
1,212 & l.~32 & \textit{for} p.~1,165 \textit{read} p.~1,175 (Lemma~1.2) \\
1,212 & l.~3 from foot & \textit{for} p.~1,290 \textit{read} p.~1,190 (Definition~2.2) \\
1,213 & l.~16 & \textit{for} stiuation \textit{read} situation \\
1,213 & l.~18 & \textit{for} infeernce \textit{read} inference \\
1,213 & l.~31 & \textit{for} Volume~IV, Chapter~6 \textit{read} Volume~IV, Chapter~5 (the claims account) \\
1,213 & l.~36 & \textit{for} conjecutre \textit{read} conjecture \\
1,213 & display (2.14) & \textit{for} $W_{j}(y) = W_{k}(z) \text{ for all } j \wedge F(y) \neq F(z)$ \textit{read} $W_{j}(y) = W_{j}(z) \text{ for all } j \wedge F(y) \neq F(z)$; the formula is used in its correct form on p.~1,217 \\
1,214 & l.~2 & \textit{for} cnovergence \textit{read} convergence \\
1,214 & l.~4 & \textit{for} demnostrate \textit{read} demonstrate \\
1,214 & l.~7 & \textit{after} isolation test \textit{insert} relative to the other six components \\
1,214 & l.~13 & \textit{for} Volume~IV, Chapter~8 \textit{read} Volume~IV, Chapter~7 (procedural fairness) \\
1,214 & l.~24 & \textit{for} p.~1,201 \textit{read} p.~1,210 (Definition~2.11) \\
1,215 & display (2.15) & \textit{for} $s_{j} \leq c_{i}$ \textit{read} $s_{j} \leq c_{j}$; correctly printed in (2.16), p.~1,216 \\
1,215 & l.~38 & \textit{for} idantify \textit{read} identify \\
1,216 & l.~25 & \textit{for} vesrion \textit{read} version \\
1,216 & l.~31 & \textit{for} diffeernce \textit{read} difference \\
1,216 & l.~39 & \textit{for} Hooker (2005) \textit{read} Broome (1990) (the claims account of fairness) \\
1,217 & l.~36 & \textit{for} capcaity \textit{read} capacity \\
1,218 & l.~4 & \textit{for} inculding \textit{read} including \\
1,218 & l.~19 & \textit{for} Volume~IV, Chapter~6 \textit{read} Volume~IV, Chapter~4 (lotteries) \\
1,218 & l.~31 & \textit{for} p.~1,270 \textit{read} p.~1,207 (Definition~2.9) \\
1,218 & l.~33 & \textit{for} essenital \textit{read} essential \\
1,218 & l.~7 from foot & \textit{for} p.~1,204 \textit{read} p.~1,203 (Theorem~2.6) \\
\end{errlist}

\subsection*{Chapter~3. Fairness and equality \hfill \normalfont\small pp.~1,219--1,247}
\begin{errlist}
1,220 & l.~1 & \textit{after} worse off \textit{insert} than others \\
1,220 & l.~8 & \textit{for} dipsute \textit{read} dispute \\
1,220 & l.~15 & \textit{for} scuceed \textit{read} succeed \\
1,220 & l.~21 & \textit{for} Definition~3.2 of Volume~I \textit{read} Definition~4.1 of Volume~I \\
1,221 & l.~8 & \textit{for} Frankfurt (1986) \textit{read} Frankfurt (1987) \\
1,221 & l.~21 & \textit{for} sentnce \textit{read} sentence \\
1,221 & l.~27 & \textit{for} Volume~IV, Chapter~7 \textit{read} Volume~IV, Chapter~9 (desert) \\
1,222 & l.~20 & \textit{delete} the second occurrence of \emph{of} \\
1,222 & l.~26 & \textit{delete} the second occurrence of \emph{a} \\
1,222 & l.~28 & \textit{for} practiacl \textit{read} practical \\
1,222 & l.~4 from foot & \textit{for} delibreate \textit{read} deliberate \\
1,223 & l.~5 & \textit{for} conssitent \textit{read} consistent \\
1,223 & l.~7 & \textit{for} $w_{k}(z) > w_{k}(x) \text{ for all } k$ \textit{read} $w_{k}(z) < w_{k}(x) \text{ for all } k$ \\
1,223 & l.~21 & \textit{for} naturlly \textit{read} naturally \\
1,223 & l.~26 & \textit{for} judgemant \textit{read} judgement \\
1,223 & l.~27 & \textit{for} p.~1,221 \textit{read} p.~1,220 (Definition~3.1) \\
1,223 & display (3.2) & \textit{for} $w_{k}(z) \leq \theta \text{ for all } k$ \textit{read} $w_{k}(z) \geq \theta \text{ for all } k$; correctly printed in (3.6), p.~1,229 \\
1,224 & l.~3 & \textit{for} obejction \textit{read} objection \\
1,224 & l.~14 & \textit{for} determiante \textit{read} determinate \\
1,225 & l.~2 & \textit{for} distinugish \textit{read} distinguish \\
1,225 & l.~4 from foot & \textit{for} p.~1,216 \textit{read} p.~1,206 (Proposition~2.7) \\
\dag\,1,225 & proof of Proposition~3.3, (3.4) & \textit{for} $w_{k}(x) \leq \theta \text{ for all } k$ \textit{read} $w_{k}(x) \geq \theta \text{ for all } k$ \\
\dag\,1,226 & proof of Theorem~3.4, l.~6 & \textit{for} Definition~3.2 \textit{read} Definition~3.1 \\
1,226 & l.~28 & \textit{for} represnt \textit{read} represent \\
\dag\,1,226 & proof of Theorem~3.4, l.~7 from foot & \textit{for} Definition~3.1 of Volume~I \textit{read} Definition~3.2 of Volume~I \\
1,226 & l.~1 from foot & \textit{for} irrleevant \textit{read} irrelevant \\
\dag\,1,227 & proof of Theorem~3.4, l.~14 & \textit{for} Definition~3.5 \textit{read} Definition~3.2; the step uses Definition~3.2 (p.~1,222) \\
1,228 & l.~14 & \textit{delete} the second occurrence of \emph{the} \\
1,228 & l.~38 & \textit{for} Volume~IV \textit{read} Volume~VI (the completeness of the list) \\
1,229 & l.~13 & \textit{for} knowldege \textit{read} knowledge \\
\dag\,1,229 & proof of Corollary~3.6, (3.5) & \textit{for} $w_{j}(z) \geq w_{j}(z) \text{ for all } j \to W(x) \geq W(z)$ \textit{read} $w_{j}(x) \geq w_{j}(z) \text{ for all } j \to W(x) \geq W(z)$ \\
1,229 & l.~20 & \textit{for} Dwokrin (1981) \textit{read} Dworkin (1981) \\
\dag\,1,229 & proof of Corollary~3.6, l.~26 & \textit{for} Definition~3.5 \textit{read} Definition~3.1; the step uses Definition~3.1 (p.~1,220) \\
1,229 & l.~29 & \textit{after} threshold \textit{insert} of sufficiency \\
1,230 & display (3.7) & \textit{for} $G(z) = \frac{1}{2n^{2}\bar{w}} \sum_{j} \sum_{k} |w_{i} - w_{k}|$ \textit{read} $G(z) = \frac{1}{2n^{2}\bar{w}} \sum_{j} \sum_{k} |w_{j} - w_{k}|$; correctly printed in (3.9), p.~1,233 \\
1,230 & l.~24 & \textit{for} the egalitarian view that a benefit matters more the worse off its recipient is in absolute terms \textit{read} the prioritarian view that a benefit matters more the worse off its recipient is in absolute terms \\
1,230 & l.~1 from foot & \textit{for} $V(x) = \sum_{j} f(w_{j}(x)), \quad f'' > 0$ \textit{read} $V(x) = \sum_{j} f(w_{j}(x)), \quad f'' < 0$ \\
1,230 & n.~5 & \textit{for} Temkin (1994) \textit{read} Temkin (1993) \\
1,231 & l.~5 & \textit{for} dsicussion \textit{read} discussion \\
1,231 & l.~9 & \textit{for} $V(x) = \sum_{i} f(w_{i}(x)), \quad f'' > 0$ \textit{read} $V(x) = \sum_{i} f(w_{i}(x)), \quad f'' < 0$ \\
1,231 & l.~22 & \textit{for} interprettaion \textit{read} interpretation \\
1,231 & display (3.8) & \textit{for} $G(z) < G(y) \wedge W(z) > W(y)$ \textit{read} $G(z) < G(y) \wedge W(z) < W(y)$ \\
1,232 & l.~27 & \textit{for} Sen (1974) \textit{read} Sen (1973) \\
1,233 & l.~3 & \textit{for} Anderson (1998) \textit{read} Anderson (1999) \\
1,233 & l.~7 & \textit{for} utlimately \textit{read} ultimately \\
1,233 & l.~17 & \textit{for} critcal \textit{read} critical \\
1,233 & l.~31 & \textit{for} riogrous \textit{read} rigorous \\
1,234 & l.~4 from foot & \textit{for} agruably \textit{read} arguably \\
1,235 & l.~17 & \textit{for} Dworkin (1980) \textit{read} Dworkin (1981) \\
1,236 & l.~32 & \textit{for} Volume~I \textit{read} Volume~II (knowledge) \\
1,237 & l.~8 & \textit{for} ambigutiy \textit{read} ambiguity \\
1,237 & l.~13 & \textit{for} $w_{j}(y) > w_{j}(x) \text{ for all } j$ \textit{read} $w_{j}(y) < w_{j}(x) \text{ for all } j$ \\
1,237 & l.~17 & \textit{delete} the second occurrence of \emph{of} \\
1,237 & l.~19 & \textit{for} Volume~IV, Chapter~5 \textit{read} Volume~IV, Chapter~6 (levelling down) \\
1,237 & l.~7 from foot & \textit{delete} the second occurrence of \emph{that} \\
1,238 & l.~9 & \textit{for} $\min_{j} w_{j}(z) > \min_{j} w_{k}(x) \to z \succ x$ \textit{read} $\min_{j} w_{j}(z) > \min_{j} w_{j}(x) \to z \succ x$ \\
1,239 & l.~3 & \textit{for} infernce \textit{read} inference \\
1,240 & l.~5 & \textit{for} Temkn (1993) \textit{read} Temkin (1993) \\
1,240 & l.~21 & \textit{for} accommdoate \textit{read} accommodate \\
1,240 & l.~4 from foot & \textit{for} respctively \textit{read} respectively \\
1,240 & n.~18 & \textit{for} Anderosn (1999) \textit{read} Anderson (1999) \\
1,241 & l.~5 & \textit{for} respnse \textit{read} response \\
1,241 & display (3.11) & \textit{for} $w_{i}(x) \geq w_{i}(x) \text{ for all } i \to W(y) \geq W(x)$ \textit{read} $w_{i}(y) \geq w_{i}(x) \text{ for all } i \to W(y) \geq W(x)$ \\
1,241 & l.~16 & \textit{for} suppsoe \textit{read} suppose \\
1,242 & l.~1 & \textit{for} indiacte \textit{read} indicate \\
1,242 & l.~5 & \textit{for} constarint \textit{read} constraint \\
\dag\,1,242 & proof of Proposition~3.7, l.~23 & \textit{for} Corollary~3.10 \textit{read} Corollary~3.6 \\
1,242 & l.~30 & \textit{for} applicaiton \textit{read} application \\
1,242 & l.~34 & \textit{for} statemnet \textit{read} statement \\
1,243 & l.~4 & \textit{delete} the second occurrence of \emph{the} \\
\dag\,1,243 & proof of Proposition~3.7, (3.13) & \textit{for} $G(y) < G(y) \wedge W(x) < W(y)$ \textit{read} $G(x) < G(y) \wedge W(x) < W(y)$ \\
1,243 & l.~6 & \textit{for} soultion \textit{read} solution \\
1,243 & l.~5 from foot & \textit{for} concerend \textit{read} concerned \\
\dag\,1,244 & proof of Lemma~3.8, (3.14) & \textit{for} $G(z) = \frac{1}{2n^{2}\bar{w}} \sum_{j} \sum_{j} |w_{j} - w_{i}|$ \textit{read} $G(z) = \frac{1}{2n^{2}\bar{w}} \sum_{j} \sum_{i} |w_{j} - w_{i}|$; the formula is used in its correct form on p.~1,247 \\
1,244 & l.~1 from foot & \textit{for} levelling up \textit{read} levelling down \\
1,245 & l.~1 & \textit{for} Volume~IV \textit{read} Volume~II (well-being) \\
1,245 & l.~6 & \textit{for} Atkinson (1969) \textit{read} Atkinson (1970) \\
1,245 & l.~8 & \textit{for} p.~1,189 \textit{read} p.~1,199 (Definition~2.3) \\
1,245 & l.~9 & \textit{for} procedurre \textit{read} procedure \\
\dag\,1,245 & proof of Lemma~3.8, l.~19 & \textit{for} Proposition~3.3 \textit{read} Proposition~3.7; the step uses Proposition~3.7 (p.~1,242) \\
\dag\,1,245 & proof of Corollary~3.10, l.~20 & \textit{for} Definition~3.5 \textit{read} Definition~3.9; the step uses Definition~3.9 (p.~1,245) \\
1,245 & l.~35 & \textit{for} valdity \textit{read} validity \\
1,246 & l.~4 & \textit{for} relative position \textit{read} absolute level \\
1,246 & l.~5 & \textit{for} Frankfurt (1988) \textit{read} Frankfurt (1987) \\
1,246 & l.~7 & \textit{for} aavilable \textit{read} available \\
1,246 & l.~28 & \textit{for} p.~1,345 \textit{read} p.~1,245 (Definition~3.9) \\
1,247 & l.~16 & \textit{for} asbtract \textit{read} abstract \\
\end{errlist}

\subsection*{Chapter~4. Lotteries and indivisible goods \hfill \normalfont\small pp.~1,248--1,283}
\begin{errlist}
1,250 & running heads, pp.~1,250--1,282 (even) & \textit{for} Fairness and equality \textit{read} Lotteries \\
1,250 & l.~29 & \textit{for} dsitinguish \textit{read} distinguish \\
1,251 & l.~3 & \textit{for} Diaomnd (1967) \textit{read} Diamond (1967) \\
1,251 & display (4.2) & \textit{for} $u_{j}(L) = \pi_{k}\, u_{k}(g)$ \textit{read} $u_{k}(L) = \pi_{k}\, u_{k}(g)$; correctly printed in (4.6), p.~1,255 \\
1,251 & l.~25 & \textit{for} $c_{i} > c_{k} \to \pi_{i} > \pi_{j}$ \textit{read} $c_{i} > c_{k} \to \pi_{i} > \pi_{k}$ \\
1,252 & l.~9 & \textit{see} Erratum~E3 \\
1,253 & l.~3 & \textit{for} dircetly \textit{read} directly \\
1,253 & display (4.5) & \textit{for} $\frac{s_{i}}{c_{j}} = \frac{s_{i}}{c_{i}}$ \textit{read} $\frac{s_{j}}{c_{j}} = \frac{s_{i}}{c_{i}}$ \\
1,253 & display (4.4) & \textit{for} $s_{i} \leq c_{j}$ \textit{read} $s_{j} \leq c_{j}$ \\
1,253 & l.~7 from foot & \textit{for} objectvie \textit{read} objective \\
1,253 & l.~6 from foot & \textit{for} p.~1,252 \textit{read} p.~1,225 (Proposition~3.3) \\
1,254 & l.~23 & \textit{for} parragraph \textit{read} paragraph \\
1,255 & l.~9 & \textit{for} opponant \textit{read} opponent \\
1,255 & l.~19 & \textit{delete} the second occurrence of \emph{is} \\
1,255 & l.~27 & \textit{for} contemporry \textit{read} contemporary \\
\dag\,1,255 & proof of Lemma~4.2, l.~30 & \textit{for} $\frac{s_{j}}{c_{i}} = \frac{s_{i}}{c_{i}}$ \textit{read} $\frac{s_{j}}{c_{j}} = \frac{s_{i}}{c_{i}}$ \\
1,256 & l.~18 & \textit{for} developmnt \textit{read} development \\
1,256 & l.~26 & \textit{for} in proportion to their urgency \textit{read} in proportion to their strength \\
1,256 & l.~31 & \textit{for} crittical \textit{read} critical \\
1,256 & l.~39 & \textit{for} Elster (1989) \textit{read} Harsanyi (1955) (the aggregation theorem) \\
1,257 & l.~5 & \textit{for} contarry \textit{read} contrary \\
\dag\,1,257 & proof of Proposition~4.3, l.~20 & \textit{for} Lemma~4.7 \textit{read} Lemma~4.2 \\
1,257 & l.~26 & \textit{for} antiicpate \textit{read} anticipate \\
1,257 & l.~30 & \textit{for} $F(x) > F(z) \wedge x_{-F} = z_{-F} \to V(x) < V(z)$ \textit{read} $F(x) > F(z) \wedge x_{-F} = z_{-F} \to V(x) > V(z)$ \\
1,257 & l.~8 from foot & \textit{for} encountr \textit{read} encounter \\
1,258 & l.~6 & \textit{for} propostion \textit{read} proposition \\
1,258 & l.~10 & \textit{for} distribtuion \textit{read} distribution \\
1,258 & display (4.8) & \textit{for} $F(x) > F(z) \wedge x_{-F} = z_{-F} \leftrightarrow V(x) > V(z)$ \textit{read} $F(x) > F(z) \wedge x_{-F} = z_{-F} \to V(x) > V(z)$ \\
1,258 & l.~18 & \textit{for} p.~1,200 \textit{read} p.~1,201 (Lemma~2.5) \\
1,258 & l.~19 & \textit{for} cosnequently \textit{read} consequently \\
1,258 & l.~20 & \textit{for} dilmma \textit{read} dilemma \\
1,258 & l.~27 & \textit{for} prnicipal \textit{read} principal \\
1,258 & l.~7 from foot & \textit{for} Stone (2010) \textit{read} Stone (2011) \\
1,258 & l.~35 & \textit{for} desrable \textit{read} desirable \\
1,260 & l.~23 & \textit{for} applictaion \textit{read} application \\
1,261 & l.~1 & \textit{for} manfest \textit{read} manifest \\
1,261 & l.~2 & \textit{for} equiavlent \textit{read} equivalent \\
1,261 & l.~26 & \textit{for} contrversy \textit{read} controversy \\
1,262 & display (4.9) & \textit{for} $\pi_{i} = c_{i} \Big/ \sum_{j} c_{i}$ \textit{read} $\pi_{i} = c_{i} \Big/ \sum_{j} c_{j}$ \\
1,262 & l.~15 & \textit{for} striclty \textit{read} strictly \\
1,262 & l.~18 & \textit{for} Wasserman (1996) \textit{read} Broome (1990) (the claims account of fairness) \\
1,262 & l.~29 & \textit{for} Elster (1990) \textit{read} Elster (1989) \\
1,263 & l.~27 & \textit{for} demontsrate \textit{read} demonstrate \\
1,263 & l.~29 & \textit{for} acknowlegde \textit{read} acknowledge \\
1,263 & l.~7 from foot & \textit{for} Eltser (1989) \textit{read} Elster (1989) \\
1,263 & l.~35 & \textit{for} Wsaserman (1996) \textit{read} Wasserman (1996) \\
1,263 & l.~3 from foot & \textit{for} impossilbe \textit{read} impossible \\
1,264 & l.~7 & \textit{for} unnceessary \textit{read} unnecessary \\
1,264 & l.~23 & \textit{for} Volume~IV \textit{read} Volume~III (relationships) \\
1,264 & l.~32 & \textit{for} altohugh \textit{read} although \\
1,264 & l.~2 from foot & \textit{for} propsect \textit{read} prospect \\
1,265 & l.~11 & \textit{for} p.~1,224 \textit{read} p.~1,242 (Proposition~3.7) \\
1,265 & l.~18 & \textit{for} Volume~IV, Chapter~10 \textit{read} Volume~IV, Chapter~9 (desert) \\
1,265 & l.~28 & \textit{for} aprpoach \textit{read} approach \\
1,266 & l.~6 & \textit{for} theerby \textit{read} thereby \\
1,266 & l.~33 & \textit{for} Sotne (2011) \textit{read} Stone (2011) \\
1,267 & l.~8 & \textit{for} partciular \textit{read} particular \\
1,267 & l.~17 & \textit{for} apporach \textit{read} approach \\
1,267 & l.~18 & \textit{for} Volume~IV, Chapter~4 \textit{read} Volume~IV, Chapter~5 (the claims account) \\
1,267 & l.~23 & \textit{for} othewrise \textit{read} otherwise \\
1,267 & l.~27 & \textit{for} abslutely \textit{read} absolutely \\
1,268 & Proposition~4.5 & \textit{see} Erratum~E2 \\
1,269 & l.~1 & \textit{for} defnitely \textit{read} definitely \\
1,269 & l.~6 from foot & \textit{for} Diamond (1968) \textit{read} Diamond (1967) \\
1,269 & n.~8 & \textit{for} Hrasanyi (1955) \textit{read} Harsanyi (1955) \\
1,270 & l.~22 & \textit{for} extenison \textit{read} extension \\
1,270 & l.~29 & \textit{for} consquence \textit{read} consequence \\
1,271 & l.~11 & \textit{for} stirctly \textit{read} strictly \\
1,271 & l.~16 & \textit{for} prcoeed \textit{read} proceed \\
1,272 & l.~9 & \textit{for} indepnedently \textit{read} independently \\
1,272 & l.~1 from foot & \textit{for} suppse \textit{read} suppose \\
1,273 & l.~15 & \textit{for} prspect \textit{read} prospect \\
1,273 & l.~21 & \textit{for} paardigm \textit{read} paradigm \\
1,273 & l.~22 & \textit{for} Harsanyi (1954) \textit{read} Harsanyi (1955) \\
1,273 & l.~25 & \textit{for} nevertheles \textit{read} nevertheless \\
1,273 & l.~4 from foot & \textit{for} alternattive \textit{read} alternative \\
1,274 & l.~14 & \textit{for} Volume~I, Chapter~6 \textit{read} Volume~I, Chapter~5 (the formal argument for plurality) \\
1,274 & l.~19 & \textit{for} Diamnod (1967) \textit{read} Diamond (1967) \\
1,274 & l.~21 & \textit{for} p.~1,286 \textit{read} p.~1,268 (Proposition~4.5) \\
1,274 & l.~22 & \textit{for} effetcively \textit{read} effectively \\
1,274 & l.~33 & \textit{for} prupose \textit{read} purpose \\
1,274 & l.~1 from foot & \textit{for} excatly \textit{read} exactly \\
1,275 & l.~24 & \textit{for} rigorrous \textit{read} rigorous \\
1,275 & l.~32 & \textit{delete} the second occurrence of \emph{the} \\
1,277 & l.~8 & \textit{for} constranit \textit{read} constraint \\
1,277 & l.~17 & \textit{for} Wassemran (1996) \textit{read} Wasserman (1996) \\
\dag\,1,277 & proof of Proposition~4.10, l.~22 & \textit{for} $c_{j} = c_{j} \to \pi_{j} = \pi_{i}$ \textit{read} $c_{j} = c_{i} \to \pi_{j} = \pi_{i}$ \\
\dag\,1,277 & proof of Proposition~4.10, (4.16) & \textit{for} $c_{j} > c_{j} \to \pi_{j} > \pi_{k}$ \textit{read} $c_{j} > c_{k} \to \pi_{j} > \pi_{k}$ \\
1,277 & l.~25 & \textit{for} eventaully \textit{read} eventually \\
\dag\,1,277 & proof of Theorem~4.8, l.~27 & \textit{for} Proposition~4.10 \textit{read} Proposition~4.5; the step uses Proposition~4.5 (p.~1,268) \\
1,278 & l.~9 & \textit{for} recongize \textit{read} recognize \\
1,278 & l.~16 & \textit{for} inevtable \textit{read} inevitable \\
\dag\,1,278 & proof of Lemma~4.11, l.~24 & \textit{for} Proposition~4.5 \textit{read} Proposition~4.3; the step uses Proposition~4.3 (p.~1,256) \\
1,278 & l.~8 from foot & \textit{for} p.~1,172 \textit{read} p.~1,272 (Theorem~4.6) \\
\dag\,1,279 & proof of Lemma~4.11, (4.17) & \textit{for} $s_{i} \geq c_{i}$ \textit{read} $s_{i} \leq c_{i}$ \\
1,279 & l.~7 from foot & \textit{for} starightforward \textit{read} straightforward \\
1,280 & l.~8 & \textit{for} ineded \textit{read} indeed \\
1,280 & l.~15 & \textit{for} Definition~3.2 of Volume~I \textit{read} Definition~3.1 of Volume~I \\
1,280 & l.~30 & \textit{for} aciheve \textit{read} achieve \\
1,280 & l.~8 from foot & \textit{for} Volume~III \textit{read} Volume~V (autonomy) \\
1,280 & l.~35 & \textit{for} consnet \textit{read} consent \\
1,281 & display (4.18) & \textit{for} $\pi_{i} = c_{i} \Big/ \sum_{k} c_{i}$ \textit{read} $\pi_{i} = c_{i} \Big/ \sum_{k} c_{k}$ \\
1,282 & l.~9 & \textit{for} ctaegory \textit{read} category \\
1,282 & l.~14 & \textit{for} cnonection \textit{read} connection \\
1,282 & l.~19 & \textit{for} ex post fairness of the lottery \textit{read} ex ante fairness of the lottery \\
1,283 & l.~9 & \textit{for} reecive \textit{read} receive \\
1,283 & l.~15 & \textit{for} p.~1,376 \textit{read} p.~1,276 (Theorem~4.8) \\
1,283 & l.~32 & \textit{for} critrion \textit{read} criterion \\
1,283 & l.~34 & \textit{for} fromulation \textit{read} formulation \\
1,283 & l.~39 & \textit{for} $\sum_{k} s_{i} \leq S$ \textit{read} $\sum_{i} s_{i} \leq S$ \\
1,283 & l.~1 from foot & \textit{for} $\sum_{j} \pi_{k} = 1$ \textit{read} $\sum_{k} \pi_{k} = 1$ \\
\end{errlist}

\subsection*{Chapter~5. Claims and their satisfaction \hfill \normalfont\small pp.~1,284--1,321}
\begin{errlist}
1,285 & l.~7 & \textit{for} irreelvant \textit{read} irrelevant \\
1,285 & l.~9 & \textit{for} rmeainder \textit{read} remainder \\
1,285 & l.~24 & \textit{for} patricular \textit{read} particular \\
1,285 & l.~35 & \textit{for} diffciulty \textit{read} difficulty \\
1,286 & l.~6 & \textit{for} unneecssary \textit{read} unnecessary \\
1,286 & l.~11 & \textit{for} chanecs \textit{read} chances \\
1,286 & l.~14 & \textit{for} surrogtae \textit{read} surrogate \\
1,286 & l.~34 & \textit{for} preimse \textit{read} premise \\
1,287 & l.~2 & \textit{for} disagrement \textit{read} disagreement \\
1,287 & l.~9 & \textit{for} adderss \textit{read} address \\
1,287 & l.~10 & \textit{for} setcion \textit{read} section \\
1,287 & l.~13 & \textit{for} Volume~IV, Chapter~3 \textit{read} Volume~IV, Chapter~4 (lotteries) \\
1,287 & l.~23 & \textit{for} Kamm (1993) \textit{read} Broome (1990) (the claims account of fairness) \\
1,287 & l.~26 & \textit{for} otbain \textit{read} obtain \\
1,287 & l.~28 & \textit{for} Broome (1992) \textit{read} Broome (1991) \\
1,288 & l.~10 & \textit{for} fairenss \textit{read} fairness \\
1,289 & l.~37 & \textit{for} serrious \textit{read} serious \\
1,290 & l.~13 & \textit{for} insetad \textit{read} instead \\
1,290 & l.~17 & \textit{for} $c_{i} > c_{k} \to \pi_{k} > \pi_{k}$ \textit{read} $c_{i} > c_{k} \to \pi_{i} > \pi_{k}$ \\
1,290 & l.~25 & \textit{for} mehtod \textit{read} method \\
1,291 & l.~8 & \textit{for} weightd \textit{read} weighted \\
1,291 & l.~26 & \textit{for} gneerally \textit{read} generally \\
1,291 & l.~34 & \textit{for} $s_{k} \leq c_{i}$ \textit{read} $s_{k} \leq c_{k}$ \\
1,292 & l.~8 & \textit{for} idetnical \textit{read} identical \\
\dag\,1,293 & proof of Proposition~5.3, l.~3 & \textit{for} $c_{k} = c_{i} \leftrightarrow \pi_{k} = \pi_{i}$ \textit{read} $c_{k} = c_{i} \to \pi_{k} = \pi_{i}$ \\
1,293 & l.~36 & \textit{for} encounetr \textit{read} encounter \\
1,294 & l.~16 & \textit{for} Volume~I, Chapter~4 \textit{read} Volume~I, Chapter~2 (the concept of the good) \\
1,295 & l.~2 & \textit{for} coheernce \textit{read} coherence \\
1,295 & l.~24 & \textit{for} p.~1,255 \textit{read} p.~1,256 (Proposition~4.3) \\
1,296 & l.~1 & \textit{for} conidtion \textit{read} condition \\
1,296 & l.~21 & \textit{for} arugably \textit{read} arguably \\
1,296 & l.~31 & \textit{for} Broome (1990) \textit{read} Hooker (2005) (fairness as a family of requirements) \\
1,296 & l.~39 & \textit{for} $u_{j}(L) = \pi_{j}\, u_{k}(g)$ \textit{read} $u_{j}(L) = \pi_{j}\, u_{j}(g)$ \\
1,297 & l.~1 & \textit{for} decisve \textit{read} decisive \\
\dag\,1,297 & proof of Lemma~5.5, (5.6) & \textit{for} $F(z) > F(x) \vee z_{-F} = x_{-F} \to V(z) > V(x)$ \textit{read} $F(z) > F(x) \wedge z_{-F} = x_{-F} \to V(z) > V(x)$; the formula is used in its correct form on p.~1,301 \\
1,297 & l.~10 & \textit{for} despitte \textit{read} despite \\
1,297 & l.~21 & \textit{for} p.~1,326 \textit{read} p.~1,226 (Theorem~3.4) \\
1,297 & l.~22 & \textit{for} impartail \textit{read} impartial \\
1,297 & l.~23 & \textit{for} remarkble \textit{read} remarkable \\
\dag\,1,298 & proof of Proposition~5.6, l.~5 & \textit{for} Lemma~5.5 \textit{read} Lemma~5.2 \\
1,298 & l.~28 & \textit{for} indvidual \textit{read} individual \\
1,299 & l.~10 & \textit{for} Kamm (1992) \textit{read} Kamm (1993) \\
1,299 & l.~11 & \textit{for} whnever \textit{read} whenever \\
\dag\,1,299 & proof of Proposition~5.6, l.~22 & \textit{for} Lemma~5.9 \textit{read} Lemma~5.5; the step uses Lemma~5.5 (p.~1,297) \\
1,299 & l.~23 & \textit{for} p.~1,298 \textit{read} p.~1,288 (Definition~5.1) \\
1,299 & l.~38 & \textit{for} chancs \textit{read} chances \\
1,300 & l.~3 & \textit{for} $c_{k} > c_{k} \to \pi_{k} > \pi_{j}$ \textit{read} $c_{k} > c_{j} \to \pi_{k} > \pi_{j}$ \\
1,300 & l.~11 & \textit{for} termnology \textit{read} terminology \\
1,301 & l.~7 & \textit{for} subsequnet \textit{read} subsequent \\
1,301 & l.~24 & \textit{for} secnario \textit{read} scenario \\
1,301 & l.~27 & \textit{for} neitehr \textit{read} neither \\
1,301 & l.~8 from foot & \textit{for} reasoanble \textit{read} reasonable \\
1,302 & l.~7 & \textit{for} scenarrio \textit{read} scenario \\
1,302 & l.~11 & \textit{for} coehrence \textit{read} coherence \\
1,302 & l.~30 & \textit{for} cotnrary \textit{read} contrary \\
1,302 & l.~33 & \textit{for} p.~1,193 \textit{read} p.~1,293 (Lemma~5.2) \\
1,303 & l.~6 & \textit{delete} the second occurrence of \emph{a} \\
1,303 & l.~18 & \textit{for} meainng \textit{read} meaning \\
1,306 & l.~12 & \textit{for} techncal \textit{read} technical \\
1,306 & l.~30 & \textit{delete} the second occurrence of \emph{that} \\
1,306 & l.~36 & \textit{for} prseence \textit{read} presence \\
1,307 & l.~2 & \textit{for} exapmle \textit{read} example \\
1,307 & l.~30 & \textit{delete} the second occurrence of \emph{is} \\
1,309 & l.~10 & \textit{for} poistion \textit{read} position \\
1,309 & l.~12 & \textit{for} procedrue \textit{read} procedure \\
1,309 & l.~13 & \textit{for} adrdess \textit{read} address \\
1,309 & l.~32 & \textit{for} p.~1,198 \textit{read} p.~1,298 (Proposition~5.6) \\
1,309 & l.~33 & \textit{for} prvided \textit{read} provided \\
1,309 & l.~5 from foot & \textit{for} pardaigm \textit{read} paradigm \\
\dag\,1,309 & proof of Lemma~5.9, (5.12) & \textit{for} $c_{j} = c_{i} \to \pi_{i} = \pi_{i}$ \textit{read} $c_{j} = c_{i} \to \pi_{j} = \pi_{i}$; the formula is used in its correct form on p.~1,310 \\
1,310 & l.~2 & \textit{for} $\frac{s_{k}}{c_{j}} = \frac{s_{j}}{c_{j}}$ \textit{read} $\frac{s_{k}}{c_{k}} = \frac{s_{j}}{c_{j}}$ \\
1,310 & l.~28 & \textit{for} regadred \textit{read} regarded \\
1,310 & l.~4 from foot & \textit{for} initilaly \textit{read} initially \\
1,311 & l.~1 & \textit{for} interst \textit{read} interest \\
\dag\,1,311 & proof of Corollary~5.11, (5.13) & \textit{for} $s_{j} \leq c_{k}$ \textit{read} $s_{j} \leq c_{j}$ \\
1,311 & l.~32 & \textit{for} evalluate \textit{read} evaluate \\
1,311 & l.~3 from foot & \textit{for} Broome (1991) \textit{read} Broome (1990) (the claims account of fairness) \\
1,312 & l.~2 from foot & \textit{for} connectoin \textit{read} connection \\
1,313 & l.~4 & \textit{for} indicte \textit{read} indicate \\
1,313 & l.~6 & \textit{for} acknolwedge \textit{read} acknowledge \\
1,313 & l.~7 & \textit{for} arugment \textit{read} argument \\
1,313 & l.~18 & \textit{for} transision \textit{read} transition \\
1,313 & l.~27 & \textit{for} $u_{k}(L) = \pi_{i}\, u_{k}(g)$ \textit{read} $u_{k}(L) = \pi_{k}\, u_{k}(g)$ \\
1,313 & l.~4 from foot & \textit{for} cateogry \textit{read} category \\
\dag\,1,314 & proof of Corollary~5.12, (5.16) & \textit{for} $\sum_{k} s_{k} \geq S$ \textit{read} $\sum_{k} s_{k} \leq S$ \\
\dag\,1,314 & proof of Corollary~5.12, l.~9 & \textit{for} Definition~5.8 \textit{read} Definition~5.10 \\
\dag\,1,314 & proof of Corollary~5.12, (5.17) & \textit{for} $F(x) > F(z) \vee x_{-F} = z_{-F} \to V(x) > V(z)$ \textit{read} $F(x) > F(z) \wedge x_{-F} = z_{-F} \to V(x) > V(z)$; correctly printed in (5.19), p.~1,315 \\
1,314 & l.~12 & \textit{for} princpal \textit{read} principal \\
1,314 & l.~34 & \textit{for} appropritae \textit{read} appropriate \\
1,314 & l.~6 from foot & \textit{for} $\pi_{j} = c_{j} \Big/ \sum_{j} c_{k}$ \textit{read} $\pi_{j} = c_{j} \Big/ \sum_{k} c_{k}$ \\
\dag\,1,314 & proof of Corollary~5.12, (5.15) & \textit{for} $\sum_{i} s_{i} \geq S$ \textit{read} $\sum_{i} s_{i} \leq S$ \\
1,315 & l.~26 & \textit{for} merley \textit{read} merely \\
1,315 & l.~8 from foot & \textit{for} questionalbe \textit{read} questionable \\
1,316 & l.~23 & \textit{for} cnosideration \textit{read} consideration \\
1,316 & l.~8 from foot & \textit{for} Brome (1991) \textit{read} Broome (1991) \\
1,317 & l.~7 & \textit{for} concernd \textit{read} concerned \\
1,317 & l.~20 & \textit{for} whehter \textit{read} whether \\
1,317 & l.~8 from foot & \textit{for} systemtic \textit{read} systematic \\
1,317 & l.~6 from foot & \textit{for} disitnction \textit{read} distinction \\
1,317 & l.~37 & \textit{for} formlation \textit{read} formulation \\
1,318 & l.~7 & \textit{for} assersion \textit{read} assertion \\
1,318 & l.~9 & \textit{for} ovbious \textit{read} obvious \\
\dag\,1,318 & proof of Proposition~5.13, l.~27 & \textit{for} Definition~5.8 \textit{read} Definition~5.7; the step uses Definition~5.7 (p.~1,300) \\
1,318 & l.~30 & \textit{for} convnetion \textit{read} convention \\
1,318 & l.~4 from foot & \textit{for} inconssitent \textit{read} inconsistent \\
1,319 & l.~12 & \textit{for} dfiference \textit{read} difference \\
1,319 & l.~15 & \textit{for} literture \textit{read} literature \\
1,319 & l.~29 & \textit{for} paragrpah \textit{read} paragraph \\
1,320 & l.~29 & \textit{for} meniton \textit{read} mention \\
1,320 & l.~6 from foot & \textit{for} evluate \textit{read} evaluate \\
1,321 & l.~8 & \textit{for} howeevr \textit{read} however \\
1,321 & l.~11 & \textit{for} intirnsic \textit{read} intrinsic \\
1,321 & l.~12 & \textit{for} p.~1,279 \textit{read} p.~1,278 (Lemma~4.11) \\
1,321 & l.~5 from foot & \textit{for} rseemblance \textit{read} resemblance \\
\end{errlist}

\subsection*{Chapter~6. Priority, sufficiency and levelling down \hfill \normalfont\small pp.~1,322--1,360}
\begin{errlist}
1,323 & l.~5 & \textit{for} exprssion \textit{read} expression \\
1,323 & l.~8 & \textit{for} $G(z) < G(x) \vee W(z) < W(x)$ \textit{read} $G(z) < G(x) \wedge W(z) < W(x)$ \\
1,323 & l.~10 & \textit{for} intendd \textit{read} intended \\
1,323 & l.~12 & \textit{for} defendd \textit{read} defended \\
1,323 & display (6.1) & \textit{for} $V(y) = \sum_{i} f(w_{i}(y)), \quad f'' > 0$ \textit{read} $V(y) = \sum_{i} f(w_{i}(y)), \quad f'' < 0$ \\
1,323 & l.~23 & \textit{for} namley \textit{read} namely \\
1,323 & l.~40 & \textit{for} consdierable \textit{read} considerable \\
1,324 & l.~17 & \textit{for} Temkin (1992) \textit{read} Temkin (1993) \\
1,324 & l.~20 & \textit{for} conditoinal \textit{read} conditional \\
1,324 & l.~25 & \textit{for} etnirely \textit{read} entirely \\
1,325 & l.~28 & \textit{for} reefrred \textit{read} referred \\
1,326 & l.~11 & \textit{for} invitable \textit{read} inevitable \\
1,326 & l.~28 & \textit{for} aboslutely \textit{read} absolutely \\
1,327 & l.~4 & \textit{for} $G(z) < G(z) \wedge W(y) < W(z)$ \textit{read} $G(y) < G(z) \wedge W(y) < W(z)$ \\
1,327 & display (6.2) & \textit{for} $V(y) = \sum_{j} f(w_{k}(y)), \quad f'' < 0$ \textit{read} $V(y) = \sum_{k} f(w_{k}(y)), \quad f'' < 0$; correctly printed in (6.4), p.~1,333 \\
1,327 & l.~15 & \textit{for} straihgtforward \textit{read} straightforward \\
1,328 & l.~20 & \textit{for} disptue \textit{read} dispute \\
1,328 & l.~23 & \textit{for} developmant \textit{read} development \\
1,329 & n.~6 & \textit{for} Crisp (2004) \textit{read} Crisp (2003) \\
1,330 & l.~15 & \textit{for} devleop \textit{read} develop \\
1,330 & l.~30 & \textit{for} obsrevation \textit{read} observation \\
1,330 & l.~8 from foot & \textit{for} preusmably \textit{read} presumably \\
1,330 & l.~5 from foot & \textit{for} deontic egalitarianism \textit{read} telic egalitarianism \\
1,331 & l.~10 & \textit{for} evidnet \textit{read} evident \\
1,331 & l.~13 & \textit{for} Holtug (2009) \textit{read} Holtug (2010) \\
1,331 & l.~8 from foot & \textit{for} euqality \textit{read} equality \\
1,332 & l.~1 & \textit{for} conseqeuntly \textit{read} consequently \\
1,332 & l.~20 & \textit{for} dependnece \textit{read} dependence \\
1,332 & l.~37 & \textit{for} preicsely \textit{read} precisely \\
1,333 & l.~3 & \textit{for} idnividual \textit{read} individual \\
1,333 & l.~14 & \textit{for} namly \textit{read} namely \\
1,333 & l.~19 & \textit{delete} the second occurrence of \emph{of} \\
\dag\,1,333 & proof of Theorem~6.3, l.~26 & \textit{for} Definition~6.2 \textit{read} Definition~6.1 \\
1,334 & l.~26 & \textit{for} setnence \textit{read} sentence \\
1,335 & l.~38 & \textit{for} contemporray \textit{read} contemporary \\
\dag\,1,336 & proof of Lemma~6.5, l.~14 & \textit{for} Definition~6.6 \textit{read} Definition~6.1 \\
1,336 & l.~18 & \textit{for} Crsp (2003) \textit{read} Crisp (2003) \\
1,337 & l.~3 & \textit{for} quessionable \textit{read} questionable \\
1,337 & l.~6 & \textit{for} conventoin \textit{read} convention \\
1,337 & l.~23 & \textit{for} calrify \textit{read} clarify \\
1,338 & l.~1 & \textit{for} immeidate \textit{read} immediate \\
1,338 & l.~17 & \textit{for} $w_{j}(y) < w_{j}(y) \text{ for all } j$ \textit{read} $w_{j}(z) < w_{j}(y) \text{ for all } j$ \\
1,338 & l.~22 & \textit{for} premse \textit{read} premise \\
1,338 & l.~31 & \textit{for} pirnciple \textit{read} principle \\
1,339 & l.~7 & \textit{for} Volume~I, Chapter~5 \textit{read} Volume~I, Chapter~4 (the method of isolation) \\
1,339 & l.~8 & \textit{for} practcial \textit{read} practical \\
1,339 & l.~23 & \textit{for} certanily \textit{read} certainly \\
1,339 & l.~37 & \textit{for} $w_{k}(y) \leq \theta \text{ for all } k$ \textit{read} $w_{k}(y) \geq \theta \text{ for all } k$ \\
1,339 & n.~22 & \textit{for} McKerlie (1995) \textit{read} McKerlie (1994) \\
1,340 & l.~1 & \textit{for} egailtarian \textit{read} egalitarian \\
1,340 & display (6.10) & \textit{for} $\min_{j} w_{j}(x) > \min_{j} w_{j}(y) \leftrightarrow x \succ y$ \textit{read} $\min_{j} w_{j}(x) > \min_{j} w_{j}(y) \to x \succ y$ \\
1,340 & l.~21 & \textit{for} frammework \textit{read} framework \\
1,340 & l.~29 & \textit{for} propossition \textit{read} proposition \\
1,341 & l.~8 from foot & \textit{for} uinform \textit{read} uniform \\
1,342 & l.~22 & \textit{for} telc \textit{read} telic \\
1,343 & l.~35 & \textit{for} determnate \textit{read} determinate \\
1,344 & l.~2 & \textit{delete} the second occurrence of \emph{the} \\
\dag\,1,344 & proof of Corollary~6.9, l.~9 & \textit{for} Definition~6.7 \textit{read} Definition~6.8 \\
1,344 & l.~10 & \textit{for} unievrsal \textit{read} universal \\
\dag\,1,344 & proof of Corollary~6.9, l.~23 & \textit{for} Definition~6.10 \textit{read} Definition~6.7; the step uses Definition~6.7 (p.~1,339) \\
1,344 & l.~28 & \textit{for} occurrecne \textit{read} occurrence \\
1,344 & l.~34 & \textit{for} existecne \textit{read} existence \\
\dag\,1,345 & proof of Corollary~6.9, (6.13) & \textit{for} $G(z) = \frac{1}{2n^{2}\bar{w}} \sum_{i} \sum_{j} |w_{i} - w_{i}|$ \textit{read} $G(z) = \frac{1}{2n^{2}\bar{w}} \sum_{i} \sum_{j} |w_{i} - w_{j}|$ \\
1,345 & l.~20 & \textit{for} questoin \textit{read} question \\
1,345 & l.~25 & \textit{for} essetnial \textit{read} essential \\
1,345 & l.~27 & \textit{for} deficeincy \textit{read} deficiency \\
1,345 & l.~8 from foot & \textit{for} appraent \textit{read} apparent \\
1,347 & l.~19 & \textit{for} McKerlie (1993) \textit{read} McKerlie (1994) \\
1,347 & display (6.14) & \textit{for} $w_{i}(x) \geq \theta \text{ for all } k$ \textit{read} $w_{k}(x) \geq \theta \text{ for all } k$ \\
1,347 & l.~30 & \textit{for} satsified \textit{read} satisfied \\
\dag\,1,348 & proof of Proposition~6.11, l.~4 & \textit{for} Definition~6.6 \textit{read} Definition~6.2 \\
\dag\,1,348 & proof of Proposition~6.11, l.~15 & \textit{for} Definition~6.10 \textit{read} Definition~6.8; the step uses Definition~6.8 (p.~1,342) \\
\dag\,1,348 & proof of Proposition~6.11, (6.16) & \textit{for} $w_{i}(x) \geq w_{i}(z) \text{ for all } i \to W(x) \leq W(z)$ \textit{read} $w_{i}(x) \geq w_{i}(z) \text{ for all } i \to W(x) \geq W(z)$ \\
\dag\,1,348 & proof of Proposition~6.11, (6.15) & \textit{for} $\min_{i} w_{i}(y) > \min_{i} w_{i}(y) \to z \succ y$ \textit{read} $\min_{i} w_{i}(z) > \min_{i} w_{i}(y) \to z \succ y$; the formula is used in its correct form on p.~1,352 \\
1,349 & l.~3 from foot & \textit{see} Erratum~E4 \\
1,350 & l.~26 & \textit{for} corrcetly \textit{read} correctly \\
\dag\,1,352 & proof of Proposition~6.12, l.~5 & \textit{for} Definition~6.6 \textit{read} Definition~6.8 \\
\dag\,1,352 & proof of Proposition~6.12, (6.20) & \textit{for} $G(y) = \frac{1}{2n^{2}\bar{w}} \sum_{k} \sum_{j} |w_{k} - w_{k}|$ \textit{read} $G(y) = \frac{1}{2n^{2}\bar{w}} \sum_{k} \sum_{j} |w_{k} - w_{j}|$ \\
1,352 & l.~15 & \textit{for} virtulaly \textit{read} virtually \\
1,352 & l.~25 & \textit{for} explanatoin \textit{read} explanation \\
1,352 & l.~32 & \textit{for} suggession \textit{read} suggestion \\
1,352 & l.~5 from foot & \textit{for} Otska and Voorhoeve (2009) \textit{read} Otsuka and Voorhoeve (2009) \\
1,353 & l.~1 & \textit{for} egalitairan \textit{read} egalitarian \\
\dag\,1,353 & proof of Proposition~6.12, l.~13 & \textit{for} Corollary~6.9 \textit{read} Corollary~6.4 \\
1,353 & l.~15 & \textit{for} circumstacne \textit{read} circumstance \\
\dag\,1,353 & proof of Proposition~6.12, l.~17 & \textit{for} Definition~6.7 \textit{read} Definition~6.2 \\
1,353 & l.~22 & \textit{for} throuhgout \textit{read} throughout \\
1,353 & l.~25 & \textit{for} obsrvation \textit{read} observation \\
1,353 & l.~1 from foot & \textit{for} implict \textit{read} implicit \\
1,354 & l.~11 & \textit{for} challnge \textit{read} challenge \\
1,354 & l.~27 & \textit{for} dilema \textit{read} dilemma \\
1,354 & l.~30 & \textit{for} implcit \textit{read} implicit \\
1,354 & l.~34 & \textit{for} p.~1,261 \textit{read} p.~1,260 (Lemma~4.4) \\
1,355 & l.~3 & \textit{for} contorversy \textit{read} controversy \\
1,355 & l.~16 & \textit{for} separte \textit{read} separate \\
1,355 & l.~25 & \textit{for} preliimnary \textit{read} preliminary \\
1,355 & l.~8 from foot & \textit{for} valdiity \textit{read} validity \\
1,356 & l.~14 & \textit{for} cnotribution \textit{read} contribution \\
1,356 & l.~30 & \textit{for} staetment \textit{read} statement \\
1,356 & n.~2 & \textit{for} Parift (1997) \textit{read} Parfit (1997) \\
1,357 & l.~5 & \textit{for} provded \textit{read} provided \\
1,357 & l.~19 & \textit{for} crtierion \textit{read} criterion \\
1,357 & l.~24 & \textit{for} p.~1,336 \textit{read} p.~1,335 (Corollary~6.4) \\
1,357 & l.~30 & \textit{for} satsfied \textit{read} satisfied \\
1,358 & l.~10 & \textit{for} proecss \textit{read} process \\
1,358 & l.~11 & \textit{for} proced \textit{read} proceed \\
1,358 & l.~32 & \textit{for} $w_{j}(z) \geq w_{k}(x) \text{ for all } j \to W(z) \geq W(x)$ \textit{read} $w_{j}(z) \geq w_{j}(x) \text{ for all } j \to W(z) \geq W(x)$ \\
1,359 & l.~7 & \textit{for} Holtg (2010) \textit{read} Holtug (2010) \\
1,359 & l.~32 & \textit{for} strngth \textit{read} strength \\
1,359 & l.~38 & \textit{for} mtoivation \textit{read} motivation \\
1,360 & l.~1 & \textit{for} a strictly convex transformation of well-being \textit{read} a strictly concave transformation of well-being \\
1,360 & l.~3 & \textit{for} intuitve \textit{read} intuitive \\
1,360 & l.~4 & \textit{for} bneeficial \textit{read} beneficial \\
1,360 & l.~9 & \textit{for} excetpion \textit{read} exception \\
1,360 & l.~22 & \textit{for} Crisp (2003) \textit{read} Frankfurt (1987) (the doctrine of sufficiency) \\
1,360 & l.~26 & \textit{for} corrseponding \textit{read} corresponding \\
1,360 & l.~30 & \textit{for} Volume~IV, Chapter~12 \textit{read} Volume~IV, Chapter~11 (the numbers) \\
1,360 & l.~37 & \textit{for} definitoin \textit{read} definition \\
\end{errlist}

\subsection*{Chapter~7. Procedural and outcome fairness \hfill \normalfont\small pp.~1,361--1,392}
\begin{errlist}
1,362 & display (7.1) & \textit{for} $c_{i} = c_{j} \to \pi_{k} = \pi_{j}$ \textit{read} $c_{i} = c_{j} \to \pi_{i} = \pi_{j}$ \\
1,362 & l.~8 from foot & \textit{for} varoius \textit{read} various \\
1,362 & l.~7 from foot & \textit{for} whetehr \textit{read} whether \\
1,362 & l.~37 & \textit{for} p.~1,313 \textit{read} p.~1,331 (Definition~6.1) \\
1,362 & l.~3 from foot & \textit{for} fundmaental \textit{read} fundamental \\
1,363 & l.~8 & \textit{for} Nozcik (1974) \textit{read} Nozick (1974) \\
1,363 & l.~12 & \textit{for} p.~1,301 \textit{read} p.~1,300 (Definition~5.7) \\
1,363 & l.~39 & \textit{for} p.~1,334 \textit{read} p.~1,344 (Corollary~6.9) \\
1,364 & l.~7 & \textit{for} oridnary \textit{read} ordinary \\
1,364 & l.~9 & \textit{for} denail \textit{read} denial \\
1,364 & l.~28 & \textit{for} cotnext \textit{read} context \\
1,364 & l.~8 from foot & \textit{for} Volume~I, Chapter~9 \textit{read} Volume~I, Chapter~8 (the seven components introduced) \\
1,365 & l.~16 & \textit{for} evidecne \textit{read} evidence \\
1,365 & l.~32 & \textit{for} moerover \textit{read} moreover \\
1,366 & l.~16 & \textit{for} conclussion \textit{read} conclusion \\
1,366 & l.~22 & \textit{for} nceessary \textit{read} necessary \\
1,366 & l.~27 & \textit{for} hypthesis \textit{read} hypothesis \\
1,366 & l.~32 & \textit{after} procedure \textit{insert} to which everyone consented \\
1,367 & l.~9 & \textit{for} pure procedural justice \textit{read} imperfect procedural justice \\
1,368 & l.~1 & \textit{for} remarkalbe \textit{read} remarkable \\
1,368 & display (7.3) & \textit{for} $u_{j}(L) = \pi_{i}\, u_{i}(g)$ \textit{read} $u_{i}(L) = \pi_{i}\, u_{i}(g)$ \\
1,368 & l.~7 & \textit{for} previuos \textit{read} previous \\
1,369 & l.~10 & \textit{for} Ralws (1958) \textit{read} Rawls (1958) \\
1,369 & l.~8 from foot & \textit{for} amdittedly \textit{read} admittedly \\
1,370 & display (7.6) & \textit{for} $c_{k} > c_{i} \to \pi_{j} > \pi_{i}$ \textit{read} $c_{j} > c_{i} \to \pi_{j} > \pi_{i}$ \\
1,370 & l.~10 & \textit{for} reasonble \textit{read} reasonable \\
1,370 & display (7.7) & \textit{for} $F(z) > F(y) \wedge z_{-F} = y_{-F} \to V(z) < V(y)$ \textit{read} $F(z) > F(y) \wedge z_{-F} = y_{-F} \to V(z) > V(y)$; correctly printed in (7.9), p.~1,373 \\
1,370 & l.~40 & \textit{for} prbable \textit{read} probable \\
1,371 & display (7.8) & \textit{for} $\frac{s_{i}}{c_{k}} = \frac{s_{k}}{c_{k}}$ \textit{read} $\frac{s_{i}}{c_{i}} = \frac{s_{k}}{c_{k}}$; the formula is used in its correct form on p.~1,372 \\
1,371 & l.~31 & \textit{for} $c_{k} < c_{j} \to \pi_{k} > \pi_{j}$ \textit{read} $c_{k} > c_{j} \to \pi_{k} > \pi_{j}$ \\
1,371 & l.~8 from foot & \textit{for} familair \textit{read} familiar \\
1,371 & l.~4 from foot & \textit{for} definision \textit{read} definition \\
1,372 & l.~14 & \textit{for} fairnses \textit{read} fairness \\
1,372 & l.~33 & \textit{for} Barry (1989) \textit{read} Hart (1955) (the principle of fair play) \\
1,372 & l.~7 from foot & \textit{for} dfeiciency \textit{read} deficiency \\
1,372 & n.~14 & \textit{for} Hart (1956) \textit{read} Hart (1955) \\
1,373 & l.~20 & \textit{for} certianly \textit{read} certainly \\
1,373 & l.~25 & \textit{for} comaprison \textit{read} comparison \\
1,374 & l.~11 & \textit{for} illusrtate \textit{read} illustrate \\
1,374 & l.~14 & \textit{for} $P(x) \wedge  P(z) \wedge x_{-F} = z_{-F}$ \textit{read} $P(x) \wedge \neg P(z) \wedge x_{-F} = z_{-F}$ \\
1,374 & l.~21 & \textit{for} clasiscal \textit{read} classical \\
1,374 & l.~34 & \textit{for} odrinary \textit{read} ordinary \\
1,374 & n.~3 & \textit{for} Rwals (1958) \textit{read} Rawls (1958) \\
1,375 & l.~10 & \textit{for} rahter \textit{read} rather \\
1,376 & l.~23 & \textit{for} obivous \textit{read} obvious \\
1,376 & l.~31 & \textit{for} empolyed \textit{read} employed \\
1,378 & l.~30 & \textit{for} anlaysis \textit{read} analysis \\
1,378 & l.~37 & \textit{for} prefernece \textit{read} preference \\
1,378 & l.~2 from foot & \textit{for} asessment \textit{read} assessment \\
\dag\,1,379 & proof of Lemma~7.5, l.~1 & \textit{for} Lemma~7.7 \textit{read} Lemma~7.2 \\
1,379 & l.~20 & \textit{for} $c_{i} = c_{j} \leftrightarrow \pi_{i} = \pi_{j}$ \textit{read} $c_{i} = c_{j} \to \pi_{i} = \pi_{j}$ \\
1,379 & l.~27 & \textit{for} permisisble \textit{read} permissible \\
\dag\,1,379 & proof of Theorem~7.3, l.~28 & \textit{for} Lemma~7.5 \textit{read} Lemma~7.2 \\
1,380 & l.~2 & \textit{for} Rawls (1971) \textit{read} Hart (1955) (the principle of fair play) \\
1,380 & l.~4 & \textit{for} acheive \textit{read} achieve \\
1,380 & l.~14 & \textit{for} p.~1,414 \textit{read} p.~1,314 (Corollary~5.12) \\
1,380 & l.~19 & \textit{for} siimlarly \textit{read} similarly \\
1,380 & l.~26 & \textit{for} Rawls (1959) \textit{read} Rawls (1958) \\
1,380 & l.~30 & \textit{for} p.~1,369 \textit{read} p.~1,379 (Lemma~7.5) \\
1,381 & l.~2 & \textit{for} rellevance \textit{read} relevance \\
1,381 & l.~17 & \textit{for} Nozick (1974) \textit{read} Hart (1955) (the principle of fair play) \\
\dag\,1,381 & proof of Lemma~7.6, (7.11) & \textit{for} $\frac{s_{k}}{c_{k}} = \frac{s_{k}}{c_{i}}$ \textit{read} $\frac{s_{k}}{c_{k}} = \frac{s_{i}}{c_{i}}$ \\
1,381 & n.~9 & \textit{for} Nzoick (1974) \textit{read} Nozick (1974) \\
1,382 & l.~6 & \textit{for} meerly \textit{read} merely \\
\dag\,1,382 & proof of Lemma~7.7, l.~13 & \textit{for} Lemma~7.6 \textit{read} Lemma~7.2; the step uses Lemma~7.2 (p.~1,372) \\
\dag\,1,382 & proof of Lemma~7.7, l.~19 & \textit{for} Lemma~7.9 \textit{read} Lemma~7.6 \\
1,382 & l.~22 & \textit{for} p.~1,381 \textit{read} p.~1,281 (Definition~4.12) \\
1,382 & l.~6 from foot & \textit{for} cnotrast \textit{read} contrast \\
1,383 & l.~4 & \textit{for} Volume~I, Chapter~7 \textit{read} Volume~I, Chapter~8 (the seven components introduced) \\
\dag\,1,383 & proof of Lemma~7.8, l.~16 & \textit{for} Lemma~7.9 \textit{read} Lemma~7.5 \\
\dag\,1,383 & proof of Lemma~7.8, l.~22 & \textit{for} Lemma~7.6 \textit{read} Lemma~7.7 \\
\dag\,1,383 & proof of Lemma~7.8, l.~26 & \textit{for} $P(x) \wedge \neg P(x) \wedge x_{-F} = y_{-F}$ \textit{read} $P(x) \wedge \neg P(y) \wedge x_{-F} = y_{-F}$ \\
1,383 & l.~30 & \textit{for} experssion \textit{read} expression \\
1,383 & l.~34 & \textit{for} particlarly \textit{read} particularly \\
1,383 & l.~35 & \textit{for} perahps \textit{read} perhaps \\
\dag\,1,384 & proof of Lemma~7.9, l.~1 & \textit{for} Lemma~7.8 \textit{read} Lemma~7.5 \\
1,384 & n.~2 & \textit{for} Nozick (1973) \textit{read} Nozick (1974) \\
1,385 & l.~10 & \textit{for} codnition \textit{read} condition \\
1,385 & l.~11 & \textit{for} p.~1,483 \textit{read} p.~1,383 (Lemma~7.8) \\
\dag\,1,385 & proof of Lemma~7.9, (7.13) & \textit{for} $\pi_{j} = c_{j} \Big/ \sum_{j} c_{i}$ \textit{read} $\pi_{j} = c_{j} \Big/ \sum_{i} c_{i}$ \\
1,385 & l.~34 & \textit{for} presupposision \textit{read} presupposition \\
1,386 & l.~11 & \textit{for} additoinal \textit{read} additional \\
1,386 & l.~16 & \textit{for} conesnt \textit{read} consent \\
1,387 & l.~2 & \textit{for} alterntaive \textit{read} alternative \\
1,387 & l.~36 & \textit{for} p.~1,322 \textit{read} p.~1,222 (Definition~3.2) \\
1,389 & l.~3 & \textit{for} p.~1,307 \textit{read} p.~1,308 (Lemma~5.9) \\
1,389 & l.~13 & \textit{for} tradittion \textit{read} tradition \\
1,390 & l.~6 & \textit{for} significnat \textit{read} significant \\
1,390 & l.~7 & \textit{for} proprotional \textit{read} proportional \\
1,390 & l.~28 & \textit{for} ocucrred \textit{read} occurred \\
1,390 & l.~4 from foot & \textit{for} oirginal \textit{read} original \\
1,391 & l.~4 & \textit{for} recirpocity \textit{read} reciprocity \\
1,391 & l.~10 & \textit{for} Volume~I, Chapter~2 \textit{read} Volume~I, Chapter~4 (the method of isolation) \\
1,391 & l.~27 & \textit{for} $\pi_{j} = c_{i} \Big/ \sum_{j} c_{j}$ \textit{read} $\pi_{i} = c_{i} \Big/ \sum_{j} c_{j}$ \\
1,391 & l.~40 & \textit{for} particulalry \textit{read} particularly \\
1,392 & l.~4 from foot & \textit{for} Brary (1989) \textit{read} Barry (1989) \\
\end{errlist}

\subsection*{Chapter~8. Fairness across generations \hfill \normalfont\small pp.~1,393--1,420}
\begin{errlist}
1,394 & l.~11 & \textit{for} expilcit \textit{read} explicit \\
1,394 & l.~7 from foot & \textit{for} accmmodate \textit{read} accommodate \\
1,395 & l.~2 & \textit{for} Parfit (1983) \textit{read} Parfit (1984) \\
1,395 & l.~26 & \textit{for} avaialble \textit{read} available \\
1,395 & l.~40 & \textit{for} p.~1,298 \textit{read} p.~1,297 (Lemma~5.5) \\
1,396 & l.~28 & \textit{for} evnetually \textit{read} eventually \\
1,396 & l.~32 & \textit{for} oppostie \textit{read} opposite \\
1,396 & l.~33 & \textit{for} precsely \textit{read} precisely \\
1,396 & l.~6 from foot & \textit{for} postrity \textit{read} posterity \\
1,396 & l.~2 from foot & \textit{for} plasuible \textit{read} plausible \\
1,397 & l.~19 & \textit{for} dennial \textit{read} denial \\
1,397 & l.~20 & \textit{for} framewrok \textit{read} framework \\
1,397 & l.~26 & \textit{for} pehraps \textit{read} perhaps \\
1,398 & l.~8 & \textit{for} trdaition \textit{read} tradition \\
1,398 & l.~17 & \textit{for} emphaiss \textit{read} emphasis \\
\dag\,1,398 & proof of Theorem~8.3, l.~27 & \textit{for} Definition~8.5 \textit{read} Definition~8.1 \\
1,398 & l.~8 from foot & \textit{after} generations \textit{insert} that do not overlap \\
1,399 & l.~14 & \textit{for} p.~1,307 \textit{read} p.~1,370 (Definition~7.1) \\
1,400 & l.~2 & \textit{for} intsead \textit{read} instead \\
1,400 & l.~7 & \textit{for} $w_{t+1}(z) \leq w_{t}(z)$ \textit{read} $w_{t+1}(z) \geq w_{t}(z)$ \\
1,400 & l.~2 from foot & \textit{for} p.~1,411 \textit{read} p.~1,311 (Definition~5.10) \\
\dag\,1,401 & proof of Proposition~8.6, l.~12 & \textit{for} Definition~8.5 \textit{read} Definition~8.4 \\
1,401 & l.~15 & \textit{for} prfeerence \textit{read} preference \\
\dag\,1,401 & proof of Proposition~8.6, l.~19 & \textit{for} Definition~8.9 \textit{read} Definition~8.5 \\
1,401 & l.~7 from foot & \textit{for} etsablish \textit{read} establish \\
1,401 & l.~5 from foot & \textit{for} Volume~VI \textit{read} Volume~V (autonomy) \\
1,402 & l.~22 & \textit{for} traetment \textit{read} treatment \\
1,402 & l.~39 & \textit{for} Parfit (1984, ch.~17) (the non-identity problem) \textit{read} Parfit (1984, ch.~16) (the non-identity problem) \\
1,403 & l.~1 & \textit{for} indeepndently \textit{read} independently \\
1,403 & l.~4 & \textit{for} wheeras \textit{read} whereas \\
1,403 & l.~7 & \textit{for} mechainsm \textit{read} mechanism \\
1,404 & l.~10 & \textit{for} requireemnt \textit{read} requirement \\
1,404 & l.~6 from foot & \textit{for} illstrate \textit{read} illustrate \\
1,405 & l.~13 & \textit{for} p.~1,294 \textit{read} p.~1,394 (Definition~8.1) \\
1,405 & l.~26 & \textit{for} presuppsition \textit{read} presupposition \\
1,405 & l.~30 & \textit{for} p.~1,343 \textit{read} p.~1,333 (Theorem~6.3) \\
\dag\,1,406 & proof of Proposition~8.8, l.~17 & \textit{for} Definition~8.9 \textit{read} Definition~8.4 \\
1,406 & l.~26 & \textit{for} comparalbe \textit{read} comparable \\
\dag\,1,406 & proof of Proposition~8.8, l.~28 & \textit{for} Proposition~8.2 \textit{read} Proposition~8.6; the step uses Proposition~8.6 (p.~1,400) \\
1,406 & l.~7 from foot & \textit{for} p.~1,326 \textit{read} p.~1,336 (Lemma~6.5) \\
1,407 & l.~1 & \textit{for} subseqeunt \textit{read} subsequent \\
1,407 & l.~4 & \textit{for} porbable \textit{read} probable \\
\dag\,1,407 & proof of Lemma~8.7, l.~14 & \textit{for} Definition~8.1 \textit{read} Definition~8.4 \\
1,407 & l.~8 from foot & \textit{for} Volume~I, Chapter~6 \textit{read} Volume~I, Chapter~7 (the refutation of monism) \\
1,408 & l.~15 & \textit{for} p.~1,240 \textit{read} p.~1,250 (Definition~4.1) \\
1,408 & l.~16 & \textit{for} immeditae \textit{read} immediate \\
1,408 & l.~23 & \textit{for} natuarlly \textit{read} naturally \\
1,408 & l.~1 from foot & \textit{for} p.~1,327 \textit{read} p.~1,317 (Proposition~5.13) \\
1,409 & l.~3 & \textit{for} p.~1,416 \textit{read} p.~1,406 (Lemma~8.7) \\
1,409 & l.~13 & \textit{for} sginificant \textit{read} significant \\
1,409 & l.~28 & \textit{for} smoething \textit{read} something \\
1,409 & l.~31 & \textit{for} a positive pure rate of time preference \textit{read} a zero pure rate of time preference \\
1,409 & l.~6 from foot & \textit{for} spcific \textit{read} specific \\
1,410 & l.~8 & \textit{for} Volume~IV \textit{read} Volume~V (beauty) \\
1,410 & l.~10 & \textit{for} unlss \textit{read} unless \\
1,410 & l.~18 & \textit{for} posetrity \textit{read} posterity \\
1,410 & l.~20 & \textit{for} appaernt \textit{read} apparent \\
1,410 & l.~23 & \textit{for} underimne \textit{read} undermine \\
1,411 & l.~9 & \textit{for} necssary \textit{read} necessary \\
1,411 & l.~12 & \textit{for} seqeunce \textit{read} sequence \\
1,411 & l.~25 & \textit{for} etxension \textit{read} extension \\
1,412 & l.~33 & \textit{for} stratgey \textit{read} strategy \\
1,412 & l.~7 from foot & \textit{for} implicaiton \textit{read} implication \\
1,413 & l.~9 & \textit{for} descirption \textit{read} description \\
1,413 & l.~26 & \textit{for} anlaogous \textit{read} analogous \\
1,414 & l.~6 & \textit{for} releavnce \textit{read} relevance \\
1,414 & l.~10 & \textit{for} resebmlance \textit{read} resemblance \\
1,414 & l.~12 & \textit{for} terimnology \textit{read} terminology \\
1,414 & l.~24 & \textit{for} undemrine \textit{read} undermine \\
1,414 & l.~25 & \textit{for} plausble \textit{read} plausible \\
1,414 & l.~29 & \textit{for} $F(x) \neq F(y) \wedge \sum_{t} w_{t}(y) = \sum_{t} w_{t}(y)$ \textit{read} $F(x) \neq F(y) \wedge \sum_{t} w_{t}(x) = \sum_{t} w_{t}(y)$ \\
1,415 & l.~18 & \textit{for} p.~1,278 \textit{read} p.~1,277 (Proposition~4.10) \\
1,415 & l.~30 & \textit{for} evidnce \textit{read} evidence \\
1,415 & l.~5 from foot & \textit{for} p.~1,500 \textit{read} p.~1,400 (Proposition~8.6) \\
1,416 & l.~8 & \textit{for} trhoughout \textit{read} throughout \\
1,417 & l.~12 & \textit{for} opopsite \textit{read} opposite \\
1,417 & l.~37 & \textit{for} attrbiute \textit{read} attribute \\
1,418 & l.~9 & \textit{for} wheenver \textit{read} whenever \\
1,419 & l.~14 & \textit{for} refreence \textit{read} reference \\
1,420 & l.~13 & \textit{for} annalogous \textit{read} analogous \\
1,420 & l.~18 & \textit{for} fruther \textit{read} further \\
1,420 & l.~27 & \textit{for} $F(y) \neq F(x) \wedge \sum_{t} w_{t}(x) = \sum_{t} w_{t}(x)$ \textit{read} $F(y) \neq F(x) \wedge \sum_{t} w_{t}(y) = \sum_{t} w_{t}(x)$ \\
1,420 & l.~31 & \textit{for} moreoevr \textit{read} moreover \\
1,420 & l.~40 & \textit{for} latetr \textit{read} latter \\
\end{errlist}

\subsection*{Chapter~9. Fairness and desert \hfill \normalfont\small pp.~1,421--1,446}
\begin{errlist}
1,422 & l.~22 & \textit{for} systemmatic \textit{read} systematic \\
1,422 & l.~25 & \textit{after} desert \textit{insert} of the recipient \\
1,422 & l.~38 & \textit{for} responsibillity \textit{read} responsibility \\
1,423 & l.~2 & \textit{for} contxt \textit{read} context \\
1,423 & l.~5 & \textit{for} Sher (1988) \textit{read} Sher (1987) \\
1,423 & l.~12 & \textit{for} phennomenon \textit{read} phenomenon \\
1,423 & l.~17 & \textit{for} oppnoent \textit{read} opponent \\
1,423 & l.~6 from foot & \textit{for} develp \textit{read} develop \\
1,423 & l.~2 from foot & \textit{for} p.~1,382 \textit{read} p.~1,372 (Lemma~7.2) \\
1,424 & l.~30 & \textit{for} logcially \textit{read} logically \\
1,425 & l.~21 & \textit{for} notasion \textit{read} notation \\
1,425 & l.~31 & \textit{for} p.~1,239 \textit{read} p.~1,339 (Definition~6.7) \\
1,426 & l.~1 & \textit{for} Miller (2000) \textit{read} Miller (1999) \\
1,426 & l.~7 & \textit{for} Volume~IV, Chapter~9 \textit{read} Volume~IV, Chapter~8 (fairness across generations) \\
1,426 & l.~9 & \textit{for} responsiiblity \textit{read} responsibility \\
1,426 & l.~23 & \textit{for} vresion \textit{read} version \\
1,426 & l.~31 & \textit{for} somehting \textit{read} something \\
1,426 & n.~2 & \textit{for} Aristotle, \emph{Nicomachean Ethics} V.4 (distributive justice) \textit{read} Aristotle, \emph{Nicomachean Ethics} V.3 (distributive justice) \\
1,427 & l.~10 & \textit{for} entrely \textit{read} entirely \\
1,428 & l.~7 & \textit{for} merrit \textit{read} merit \\
1,428 & l.~10 & \textit{for} Feinberg (1969) \textit{read} Feinberg (1970) \\
1,428 & l.~32 & \textit{for} acocunt \textit{read} account \\
1,428 & l.~8 from foot & \textit{for} clamied \textit{read} claimed \\
1,429 & l.~10 & \textit{for} Fienberg (1970) \textit{read} Feinberg (1970) \\
1,429 & l.~13 & \textit{for} Volume~I, Chapter~9 \textit{read} Volume~I, Chapter~7 (the refutation of monism) \\
1,429 & l.~15 & \textit{for} circumtsance \textit{read} circumstance \\
1,429 & l.~38 & \textit{for} p.~1,397 \textit{read} p.~1,398 (Theorem~8.3) \\
1,430 & l.~27 & \textit{for} copmarison \textit{read} comparison \\
1,430 & l.~37 & \textit{for} anticpiate \textit{read} anticipate \\
\dag\,1,431 & proof of Proposition~9.5, l.~1 & \textit{for} Proposition~9.7 \textit{read} Proposition~9.4; the step uses Proposition~9.4 (p.~1,431) \\
1,431 & l.~20 & \textit{for} nothnig \textit{read} nothing \\
1,431 & l.~22 & \textit{for} intrest \textit{read} interest \\
\dag\,1,432 & proof of Proposition~9.5, l.~3 & \textit{for} Proposition~9.7 \textit{read} Proposition~9.2 \\
\dag\,1,432 & proof of Proposition~9.4, l.~17 & \textit{for} Proposition~9.5 \textit{read} Proposition~9.2; the step uses Proposition~9.2 (p.~1,424) \\
1,432 & l.~21 & \textit{for} idnetify \textit{read} identify \\
1,432 & l.~5 from foot & \textit{for} p.~1,396 \textit{read} p.~1,397 (Proposition~8.2) \\
1,432 & l.~1 from foot & \textit{for} separtae \textit{read} separate \\
1,433 & l.~38 & \textit{for} desert-insensitive \textit{read} desert-sensitive \\
1,433 & n.~20 & \textit{for} Miller (1998) \textit{read} Miller (1999) \\
\dag\,1,434 & proof of Proposition~9.7, (9.5) & \textit{for} $F(x) = -\sum_{j} |s_{i} - d_{j}|$ \textit{read} $F(x) = -\sum_{j} |s_{j} - d_{j}|$ \\
\dag\,1,434 & proof of Proposition~9.7, l.~29 & \textit{for} Proposition~9.2 \textit{read} Proposition~9.4; the step uses Proposition~9.4 (p.~1,431) \\
\dag\,1,435 & proof of Proposition~9.7, l.~13 & \textit{for} Proposition~9.4 \textit{read} Proposition~9.5; the step uses Proposition~9.5 (p.~1,431) \\
1,435 & l.~20 & \textit{for} remaindr \textit{read} remainder \\
1,436 & l.~5 & \textit{for} ulness \textit{read} unless \\
1,436 & l.~18 & \textit{for} comlpete \textit{read} complete \\
1,437 & n.~13 & \textit{for} Kagan (2011) \textit{read} Kagan (2012) \\
1,438 & l.~3 & \textit{for} acconut \textit{read} account \\
1,438 & l.~18 & \textit{for} deesrt \textit{read} desert \\
1,438 & l.~23 & \textit{for} $d_{k} > d_{i} \to s_{k} \leq s_{i}$ \textit{read} $d_{k} > d_{i} \to s_{k} \geq s_{i}$ \\
1,438 & l.~27 & \textit{for} p.~1,499 \textit{read} p.~1,399 (Definition~8.5) \\
1,438 & l.~4 from foot & \textit{for} wheraes \textit{read} whereas \\
1,439 & l.~2 & \textit{for} seriuos \textit{read} serious \\
1,439 & l.~7 & \textit{for} challnege \textit{read} challenge \\
1,439 & l.~9 & \textit{for} itnuitive \textit{read} intuitive \\
1,439 & l.~17 & \textit{for} concluison \textit{read} conclusion \\
1,439 & l.~30 & \textit{for} indepednent \textit{read} independent \\
1,440 & l.~2 & \textit{for} Kgaan (2012) \textit{read} Kagan (2012) \\
1,440 & display (9.8) & \textit{for} $d_{j} > d_{k} \to s_{j} \geq s_{i}$ \textit{read} $d_{j} > d_{k} \to s_{j} \geq s_{k}$; correctly printed in (9.12), p.~1,442 \\
1,440 & l.~22 & \textit{for} comparative desert \textit{read} noncomparative desert \\
1,440 & display (9.9) & \textit{for} $F(x) = -\sum_{i} |s_{j} - d_{j}|$ \textit{read} $F(x) = -\sum_{j} |s_{j} - d_{j}|$ \\
1,440 & l.~4 from foot & \textit{for} Feinberg (1971) \textit{read} Feinberg (1970) \\
1,441 & l.~3 & \textit{for} proposl \textit{read} proposal \\
1,441 & l.~7 & \textit{for} therrefore \textit{read} therefore \\
1,441 & l.~10 & \textit{for} nevetrheless \textit{read} nevertheless \\
1,441 & display (9.10) & \textit{for} $s_{j} / d_{k} = s_{i} / d_{i}$ \textit{read} $s_{j} / d_{j} = s_{i} / d_{i}$; correctly printed in (9.11), p.~1,442 \\
1,441 & l.~21 & \textit{for} initialy \textit{read} initially \\
1,441 & l.~5 from foot & \textit{for} p.~1,433 \textit{read} p.~1,423 (Definition~9.1) \\
1,442 & l.~40 & \textit{for} p.~1,432 \textit{read} p.~1,332 (Definition~6.2) \\
1,443 & display (9.13) & \textit{for} $s_{i} / d_{i} = s_{i} / d_{k}$ \textit{read} $s_{i} / d_{i} = s_{k} / d_{k}$; the formula is used in its correct form on p.~1,444 \\
1,444 & l.~18 & \textit{for} assumpiton \textit{read} assumption \\
1,444 & l.~28 & \textit{for} threefore \textit{read} therefore \\
1,445 & l.~26 & \textit{for} p.~1,479 \textit{read} p.~1,379 (Proposition~7.4) \\
1,445 & l.~1 from foot & \textit{for} p.~1,411 \textit{read} p.~1,311 (Corollary~5.11) \\
1,446 & l.~5 & \textit{for} gennuine \textit{read} genuine \\
1,446 & l.~8 & \textit{for} comparbale \textit{read} comparable \\
1,446 & l.~15 & \textit{for} p.~1,346 \textit{read} p.~1,436 (Corollary~9.8) \\
1,446 & l.~32 & \textit{for} Kaagn (2012) \textit{read} Kagan (2012) \\
1,446 & l.~3 from foot & \textit{for} later \textit{read} latter \\
\end{errlist}

\subsection*{Chapter~10. Fairness is not a component of well-being \hfill \normalfont\small pp.~1,447--1,477}
\begin{errlist}
1,448 & l.~6 & \textit{for} dependnce \textit{read} dependence \\
1,448 & l.~7 & \textit{for} $G(x) < G(z) \wedge W(z) < W(z)$ \textit{read} $G(x) < G(z) \wedge W(x) < W(z)$ \\
1,448 & l.~27 & \textit{for} Arnesn (1989) \textit{read} Arneson (1989) \\
1,449 & l.~26 & \textit{for} $\min_{k} w_{k}(z) < \min_{k} w_{k}(x) \to z \succ x$ \textit{read} $\min_{k} w_{k}(z) > \min_{k} w_{k}(x) \to z \succ x$ \\
1,450 & l.~3 & \textit{for} equlaly \textit{read} equally \\
1,450 & l.~14 & \textit{for} thershold \textit{read} threshold \\
1,451 & l.~3 & \textit{for} decsribe \textit{read} describe \\
1,451 & l.~6 & \textit{for} Hirose (2015) \textit{read} Cohen (1989) (access to advantage) \\
1,451 & l.~12 & \textit{for} p.~1,454 \textit{read} p.~1,445 (Definition~9.9) \\
1,451 & l.~25 & \textit{for} conesnsus \textit{read} consensus \\
1,451 & l.~26 & \textit{for} equaly \textit{read} equally \\
1,451 & display (10.1) & \textit{for} $w_{j}(z) < w_{j}(z) \text{ for all } j$ \textit{read} $w_{j}(x) < w_{j}(z) \text{ for all } j$ \\
1,451 & l.~38 & \textit{for} $y_{-F} = x_{-F} \wedge y_{F} \neq y_{F} \wedge V(y) \neq V(x)$ \textit{read} $y_{-F} = x_{-F} \wedge y_{F} \neq x_{F} \wedge V(y) \neq V(x)$ \\
1,452 & l.~9 & \textit{for} cmomitment \textit{read} commitment \\
1,452 & l.~16 & \textit{for} emphsais \textit{read} emphasis \\
1,452 & l.~28 & \textit{for} varitaion \textit{read} variation \\
1,452 & display (10.2) & \textit{for} $w_{i}(z) \leq \theta \text{ for all } i$ \textit{read} $w_{i}(z) \geq \theta \text{ for all } i$ \\
1,452 & l.~5 from foot & \textit{for} inquality \textit{read} inequality \\
1,453 & l.~19 & \textit{for} $w_{j}(z) \leq w_{j}(y) \text{ for all } j \to W(z) \geq W(y)$ \textit{read} $w_{j}(z) \geq w_{j}(y) \text{ for all } j \to W(z) \geq W(y)$ \\
1,453 & display (10.3) & \textit{for} $\min_{j} w_{j}(y) < \min_{j} w_{j}(z) \to y \succ z$ \textit{read} $\min_{j} w_{j}(y) > \min_{j} w_{j}(z) \to y \succ z$ \\
1,453 & l.~31 & \textit{for} uniofrm \textit{read} uniform \\
1,453 & l.~33 & \textit{for} carfeully \textit{read} carefully \\
1,453 & l.~7 from foot & \textit{for} carefuly \textit{read} carefully \\
1,453 & l.~4 from foot & \textit{for} stratgy \textit{read} strategy \\
1,454 & l.~9 & \textit{for} adequtae \textit{read} adequate \\
1,454 & l.~11 & \textit{for} Paretto \textit{read} Pareto \\
1,455 & display (10.4) & \textit{for} $w_{j}(x) \geq w_{j}(y) \text{ for all } j \to W(y) \geq W(y)$ \textit{read} $w_{j}(x) \geq w_{j}(y) \text{ for all } j \to W(x) \geq W(y)$; the formula is used in its correct form on p.~1,458 \\
1,455 & l.~40 & \textit{for} estalbish \textit{read} establish \\
1,456 & l.~15 & \textit{for} monsim \textit{read} monism \\
1,456 & l.~31 & \textit{for} maening \textit{read} meaning \\
1,457 & l.~20 & \textit{for} $w_{i}(y) > w_{i}(x) \text{ for all } i$ \textit{read} $w_{i}(y) < w_{i}(x) \text{ for all } i$ \\
1,457 & l.~2 from foot & \textit{for} thoerem \textit{read} theorem \\
1,458 & l.~3 from foot & \textit{for} Arneson (1989) \textit{read} Cohen (1989) (access to advantage) \\
1,459 & proof of Theorem~10.4 & \textit{see} Erratum~E5 \\
1,460 & l.~2 & \textit{for} emplyoed \textit{read} employed \\
1,461 & l.~23 & \textit{for} occrured \textit{read} occurred \\
1,462 & l.~8 & \textit{for} calimed \textit{read} claimed \\
1,462 & l.~15 & \textit{for} consesnus \textit{read} consensus \\
1,462 & l.~36 & \textit{for} p.~1,470 \textit{read} p.~1,460 (Definition~10.6) \\
1,462 & l.~37 & \textit{for} Scanlon (1999) \textit{read} Scanlon (1998) \\
1,462 & l.~39 & \textit{for} explciit \textit{read} explicit \\
1,463 & l.~2 & \textit{for} p.~1,468 \textit{read} p.~1,458 (Theorem~10.2) \\
1,463 & l.~3 & \textit{for} hypothseis \textit{read} hypothesis \\
1,463 & l.~18 & \textit{for} suggetsion \textit{read} suggestion \\
1,463 & display (10.5) & \textit{for} $x \succ y \iff V(x) < V(y)$ \textit{read} $x \succ y \iff V(x) > V(y)$ \\
1,463 & l.~1 from foot & \textit{for} empriical \textit{read} empirical \\
\dag\,1,464 & proof of Lemma~10.7, l.~6 & \textit{for} Definition~10.6 \textit{read} Definition~10.1 \\
\dag\,1,464 & proof of Lemma~10.7, (10.6) & \textit{for} $\forall z\,\forall z\,(x|_{S} = z|_{S} \to V(x) = V(z))$ \textit{read} $\forall x\,\forall z\,(x|_{S} = z|_{S} \to V(x) = V(z))$ \\
\dag\,1,464 & proof of Lemma~10.7, l.~11 & \textit{for} Theorem~10.2 \textit{read} Theorem~10.3; the step uses Theorem~10.3 (p.~1,459) \\
1,464 & l.~19 & \textit{for} concvae \textit{read} concave \\
1,464 & l.~34 & \textit{for} cocneption \textit{read} conception \\
1,464 & l.~2 from foot & \textit{for} lanugage \textit{read} language \\
1,465 & l.~9 & \textit{for} identcial \textit{read} identical \\
\dag\,1,465 & proof of Lemma~10.8, (10.7) & \textit{for} $G(y) = \frac{1}{2n^{2}\bar{w}} \sum_{i} \sum_{j} |w_{j} - w_{j}|$ \textit{read} $G(y) = \frac{1}{2n^{2}\bar{w}} \sum_{i} \sum_{j} |w_{i} - w_{j}|$; correctly printed in (10.12), p.~1,470 \\
\dag\,1,465 & proof of Lemma~10.8, l.~24 & \textit{for} Theorem~10.4 \textit{read} Theorem~10.3 \\
\dag\,1,465 & proof of Lemma~10.8, l.~28 & \textit{for} Lemma~10.10 \textit{read} Lemma~10.7; the step uses Lemma~10.7 (p.~1,464) \\
1,465 & n.~21 & \textit{for} Choen (1989) \textit{read} Cohen (1989) \\
1,466 & l.~10 & \textit{for} cnostruction \textit{read} construction \\
1,466 & l.~15 & \textit{after} idealized reflective endorsement \textit{insert} from the impartial standpoint \\
\dag\,1,466 & proof of Corollary~10.9, (10.8) & \textit{for} $w_{k}(y) \geq \theta \text{ for all } j$ \textit{read} $w_{j}(y) \geq \theta \text{ for all } j$ \\
\dag\,1,467 & proof of Lemma~10.10, (10.9) & \textit{for} $\forall y\,\forall y\,(y|_{S} = z|_{S} \to V(y) = V(z))$ \textit{read} $\forall y\,\forall z\,(y|_{S} = z|_{S} \to V(y) = V(z))$ \\
1,467 & l.~5 & \textit{for} geniune \textit{read} genuine \\
\dag\,1,467 & proof of Corollary~10.9, l.~10 & \textit{for} Lemma~10.8 \textit{read} Lemma~10.7; the step uses Lemma~10.7 (p.~1,464) \\
1,467 & l.~19 & \textit{for} otucome \textit{read} outcome \\
1,467 & l.~25 & \textit{for} mnoism \textit{read} monism \\
1,469 & display (10.11) & \textit{for} $z|_{S} = y|_{S} \vee V(z) \neq V(y)$ \textit{read} $z|_{S} = y|_{S} \wedge V(z) \neq V(y)$; correctly printed in (10.13), p.~1,470 \\
1,469 & l.~3 & \textit{for} Scanoln (1998) \textit{read} Scanlon (1998) \\
1,469 & l.~37 & \textit{for} Cohen (1988) \textit{read} Cohen (1989) \\
1,469 & display (10.10) & \textit{for} $G(z) = \frac{1}{2n^{2}\bar{w}} \sum_{k} \sum_{i} |w_{i} - w_{i}|$ \textit{read} $G(z) = \frac{1}{2n^{2}\bar{w}} \sum_{k} \sum_{i} |w_{k} - w_{i}|$ \\
1,470 & l.~10 & \textit{for} omited \textit{read} omitted \\
1,470 & n.~24 & \textit{for} Cohn (1989) \textit{read} Cohen (1989) \\
1,471 & display (10.14) & \textit{for} $W_{k}(y) = W_{k}(z) \text{ for all } k \vee F(y) \neq F(z)$ \textit{read} $W_{k}(y) = W_{k}(z) \text{ for all } k \wedge F(y) \neq F(z)$ \\
1,471 & l.~1 from foot & \textit{for} Hirose (2014) \textit{read} Hirose (2015) \\
1,472 & l.~1 & \textit{for} $V(z) = \sum_{k} f(w_{i}(z)), \quad f'' < 0$ \textit{read} $V(z) = \sum_{k} f(w_{k}(z)), \quad f'' < 0$ \\
1,472 & l.~22 & \textit{for} Arneson (1988) \textit{read} Arneson (1989) \\
1,472 & display (10.15) & \textit{for} $y|_{S} = z|_{S} \wedge V(y) \neq V(y)$ \textit{read} $y|_{S} = z|_{S} \wedge V(y) \neq V(z)$ \\
1,472 & l.~32 & \textit{for} p.~1,425 \textit{read} p.~1,245 (Corollary~3.10) \\
1,472 & l.~4 from foot & \textit{for} assumtpion \textit{read} assumption \\
1,472 & l.~2 from foot & \textit{for} ienquality \textit{read} inequality \\
1,473 & l.~21 & \textit{for} $y_{-F} = x_{-F} \wedge x_{F} \neq x_{F} \wedge V(y) \neq V(x)$ \textit{read} $y_{-F} = x_{-F} \wedge y_{F} \neq x_{F} \wedge V(y) \neq V(x)$ \\
1,473 & l.~23 & \textit{for} pehnomenon \textit{read} phenomenon \\
1,473 & l.~32 & \textit{for} rseponse \textit{read} response \\
1,474 & l.~2 & \textit{for} equlaity \textit{read} equality \\
1,474 & display (10.16) & \textit{for} $x \succ z \iff V(x) > V(x)$ \textit{read} $x \succ z \iff V(x) > V(z)$ \\
1,474 & l.~12 & \textit{for} rfelective \textit{read} reflective \\
1,475 & l.~5 & \textit{for} prinicple \textit{read} principle \\
1,475 & l.~11 & \textit{for} Cohen (1989) \textit{read} Arneson (1989) (equal opportunity for welfare) \\
1,475 & l.~14 & \textit{for} Parreto \textit{read} Pareto \\
1,475 & l.~2 from foot & \textit{for} dificulty \textit{read} difficulty \\
1,476 & l.~6 & \textit{for} smiilarly \textit{read} similarly \\
1,476 & l.~24 & \textit{for} profle \textit{read} profile \\
1,477 & l.~20 & \textit{for} adqeuate \textit{read} adequate \\
1,477 & l.~24 & \textit{for} treatmant \textit{read} treatment \\
1,477 & l.~28 & \textit{for} p.~1,405 \textit{read} p.~1,406 (Proposition~8.8) \\
\end{errlist}

\subsection*{Chapter~11. Aggregation and the numbers \hfill \normalfont\small pp.~1,478--1,507}
\begin{errlist}
1,479 & l.~6 & \textit{for} p.~1,238 \textit{read} p.~1,228 (Definition~3.5) \\
1,479 & l.~25 & \textit{for} Otsuka and Voorhoeve (2009) \textit{read} Taurek (1977) (the numbers problem) \\
1,480 & l.~2 & \textit{for} the numbers do not count \textit{read} the numbers count \\
1,480 & l.~3 & \textit{for} instacne \textit{read} instance \\
1,480 & l.~9 & \textit{for} undertsanding \textit{read} understanding \\
1,480 & l.~14 & \textit{for} varaition \textit{read} variation \\
1,480 & l.~30 & \textit{for} p.~1,265 \textit{read} p.~1,255 (Lemma~4.2) \\
1,480 & l.~4 from foot & \textit{for} understnading \textit{read} understanding \\
1,481 & l.~11 & \textit{for} apropriate \textit{read} appropriate \\
1,481 & l.~19 & \textit{for} disagreeemnt \textit{read} disagreement \\
1,481 & l.~35 & \textit{for} aggergation \textit{read} aggregation \\
1,482 & l.~1 & \textit{for} manitain \textit{read} maintain \\
1,482 & l.~8 & \textit{for} p.~1,414 \textit{read} p.~1,424 (Proposition~9.2) \\
1,482 & l.~10 & \textit{for} the tie-breaking argument assumes aggregation \textit{read} the tie-breaking argument does not assume aggregation \\
1,482 & l.~20 & \textit{for} descirbe \textit{read} describe \\
1,482 & l.~27 & \textit{for} p.~1,318 \textit{read} p.~1,381 (Lemma~7.6) \\
1,483 & l.~1 & \textit{for} p.~1,328 \textit{read} p.~1,338 (Definition~6.6) \\
1,483 & l.~34 & \textit{for} p.~1,279 \textit{read} p.~1,379 (Theorem~7.3) \\
1,483 & l.~5 from foot & \textit{for} contractaulist \textit{read} contractualist \\
1,484 & l.~4 & \textit{for} oriignal \textit{read} original \\
1,484 & l.~17 & \textit{for} refrered \textit{read} referred \\
\dag\,1,484 & proof of Proposition~11.3, l.~24 & \textit{for} Definition~11.4 \textit{read} Definition~11.1; the step uses Definition~11.1 (p.~1,482) \\
1,484 & l.~27 & \textit{for} p.~1,463 \textit{read} p.~1,464 (Lemma~10.7) \\
1,485 & l.~3 & \textit{for} omitetd \textit{read} omitted \\
1,485 & l.~29 & \textit{for} srtength \textit{read} strength \\
1,485 & l.~31 & \textit{for} propponent \textit{read} proponent \\
1,486 & l.~4 & \textit{for} indviidualist \textit{read} individualist \\
1,487 & l.~25 & \textit{for} p.~1,432 \textit{read} p.~1,431 (Proposition~9.4) \\
\dag\,1,488 & proof of Lemma~11.7, l.~8 & \textit{for} Lemma~11.5 \textit{read} Lemma~11.6; the step uses Lemma~11.6 (p.~1,487) \\
\dag\,1,488 & proof of Lemma~11.7, l.~11 & \textit{for} Definition~11.8 \textit{read} Definition~11.4 \\
1,488 & l.~20 & \textit{for} p.~1,282 \textit{read} p.~1,382 (Lemma~7.7) \\
\dag\,1,488 & proof of Lemma~11.6, l.~24 & \textit{for} Lemma~11.10 \textit{read} Lemma~11.5 \\
1,488 & l.~28 & \textit{for} p.~1,483 \textit{read} p.~1,482 (Definition~11.1) \\
1,489 & l.~7 & \textit{for} corrctly \textit{read} correctly \\
1,489 & l.~10 & \textit{for} p.~1,477 \textit{read} p.~1,487 (Lemma~11.6) \\
1,489 & l.~12 & \textit{for} p.~1,309 \textit{read} p.~1,409 (Definition~8.9) \\
1,489 & l.~19 & \textit{for} individaulist \textit{read} individualist \\
1,489 & l.~26 & \textit{for} $n_{A} > n_{B} \leftrightarrow A \succ B$ \textit{read} $n_{A} > n_{B} \to A \succ B$ \\
1,490 & l.~17 & \textit{for} betwen \textit{read} between \\
1,490 & l.~26 & \textit{for} perusasive \textit{read} persuasive \\
1,491 & l.~4 & \textit{for} p.~1,325 \textit{read} p.~1,352 (Proposition~6.12) \\
1,491 & l.~5 & \textit{for} p.~1,478 \textit{read} p.~1,488 (Lemma~11.7) \\
1,491 & l.~23 & \textit{for} commitmnet \textit{read} commitment \\
1,491 & l.~27 & \textit{for} mantion \textit{read} mention \\
1,492 & l.~15 & \textit{for} p.~1,456 \textit{read} p.~1,465 (Lemma~10.8) \\
1,492 & l.~31 & \textit{for} fundmental \textit{read} fundamental \\
1,493 & l.~22 & \textit{for} contsruction \textit{read} construction \\
1,493 & l.~27 & \textit{for} p.~1,413 \textit{read} p.~1,431 (Proposition~9.5) \\
1,494 & l.~9 & \textit{for} structre \textit{read} structure \\
1,494 & l.~4 from foot & \textit{for} p.~1,328 \textit{read} p.~1,228 (Corollary~3.6) \\
1,495 & l.~8 & \textit{for} intsance \textit{read} instance \\
1,495 & l.~18 & \textit{for} conjecturre \textit{read} conjecture \\
1,495 & l.~25 & \textit{for} pairwse \textit{read} pairwise \\
1,495 & l.~26 & \textit{for} preesnce \textit{read} presence \\
1,495 & l.~28 & \textit{for} pairiwse \textit{read} pairwise \\
1,496 & l.~1 & \textit{for} Kamm (1993) \textit{read} Otsuka and Voorhoeve (2009) (the argument against the priority view) \\
1,496 & l.~3 & \textit{for} indepndent \textit{read} independent \\
1,496 & l.~34 & \textit{for} Taurek (1976) \textit{read} Taurek (1977) \\
1,497 & l.~15 & \textit{after} pairwise comparison \textit{insert} of individual claims \\
1,498 & l.~11 & \textit{for} otherwse \textit{read} otherwise \\
1,499 & l.~10 & \textit{for} p.~1,285 \textit{read} p.~1,385 (Definition~7.10) \\
1,499 & l.~14 & \textit{for} p.~1,277 \textit{read} p.~1,276 (Corollary~4.9) \\
1,499 & l.~16 & \textit{for} empiriacl \textit{read} empirical \\
1,499 & l.~31 & \textit{for} redudnancy \textit{read} redundancy \\
1,500 & l.~1 & \textit{for} p.~1,299 \textit{read} p.~1,399 (Definition~8.4) \\
1,500 & l.~24 & \textit{for} convction \textit{read} conviction \\
1,500 & l.~2 from foot & \textit{for} defiintely \textit{read} definitely \\
1,501 & l.~31 & \textit{for} p.~1,375 \textit{read} p.~1,275 (Lemma~4.7) \\
1,502 & l.~2 & \textit{for} Kamm (1993) \textit{read} Taurek (1977) (the numbers problem) \\
1,502 & l.~4 & \textit{for} p.~1,548 \textit{read} p.~1,448 (Definition~10.1) \\
1,502 & l.~13 & \textit{for} propoennt \textit{read} proponent \\
1,502 & l.~14 & \textit{for} p.~1,491 \textit{read} p.~1,490 (Definition~11.9) \\
1,503 & l.~2 & \textit{for} situatoin \textit{read} situation \\
1,504 & l.~8 & \textit{for} includng \textit{read} including \\
1,504 & l.~18 & \textit{for} p.~1,458 \textit{read} p.~1,459 (Theorem~10.4) \\
1,504 & l.~33 & \textit{for} atrtibute \textit{read} attribute \\
1,505 & l.~25 & \textit{for} exctly \textit{read} exactly \\
1,505 & l.~26 & \textit{for} reagrdless \textit{read} regardless \\
1,505 & l.~29 & \textit{for} mechansim \textit{read} mechanism \\
1,505 & l.~35 & \textit{for} beneifcial \textit{read} beneficial \\
1,505 & n.~23 & \textit{for} Taurk (1977) \textit{read} Taurek (1977) \\
\dag\,1,506 & proof of Lemma~11.10, l.~1 & \textit{for} Definition~11.8 \textit{read} Definition~11.9 \\
\dag\,1,506 & proof of Lemma~11.10, l.~23 & \textit{for} Lemma~11.7 \textit{read} Lemma~11.6; the step uses Lemma~11.6 (p.~1,487) \\
\dag\,1,506 & proof of Lemma~11.10, l.~28 & \textit{for} Lemma~11.7 \textit{read} Lemma~11.5 \\
\dag\,1,507 & proof of Lemma~11.10, l.~2 & \textit{for} Definition~11.9 \textit{read} Definition~11.4; the step uses Definition~11.4 (p.~1,484) \\
\dag\,1,507 & proof of Lemma~11.10, l.~5 & \textit{for} $n_{A} < n_{B} \to A \succ B$ \textit{read} $n_{A} > n_{B} \to A \succ B$ \\
1,507 & l.~11 & \textit{for} aggregtaion \textit{read} aggregation \\
1,507 & l.~28 & \textit{for} Scanlon (1998) \textit{read} Taurek (1977) (the numbers problem) \\
1,507 & l.~30 & \textit{for} p.~1,505 \textit{read} p.~1,506 (Lemma~11.10) \\
1,507 & l.~4 from foot & \textit{for} Volume~VI \textit{read} Volume~V (beauty) \\
\end{errlist}

\subsection*{Chapter~12. Summary: fairness in the list of goods \hfill \normalfont\small pp.~1,508--1,520}
\begin{errlist}
1,509 & l.~3 & \textit{for} irreducbile \textit{read} irreducible \\
1,509 & l.~10 & \textit{for} annalysis \textit{read} analysis \\
1,509 & l.~17 & \textit{for} restritcion \textit{read} restriction \\
1,509 & l.~29 & \textit{for} Rawls (1971) \textit{read} Broome (1990) (the claims account of fairness) \\
1,509 & l.~30 & \textit{for} occasionaly \textit{read} occasionally \\
1,510 & running heads, pp.~1,510--1,520 (even) & \textit{for} Aggregation and the numbers \textit{read} Summary \\
1,510 & l.~9 & \textit{for} p.~1,473 \textit{read} p.~1,483 (Proposition~11.3) \\
1,510 & l.~23 & \textit{for} endorsemnt \textit{read} endorsement \\
1,510 & l.~24 & \textit{for} speicfic \textit{read} specific \\
1,510 & l.~25 & \textit{for} scetion \textit{read} section \\
1,510 & l.~1 from foot & \textit{for} porfile \textit{read} profile \\
1,511 & l.~1 & \textit{for} accordngly \textit{read} accordingly \\
1,511 & l.~3 & \textit{for} p.~1,432 \textit{read} p.~1,342 (Definition~6.8) \\
1,511 & l.~5 & \textit{for} acocrdingly \textit{read} accordingly \\
1,511 & l.~18 & \textit{for} occasionlaly \textit{read} occasionally \\
1,511 & l.~27 & \textit{for} volme \textit{read} volume \\
1,511 & l.~32 & \textit{for} idealiezd \textit{read} idealized \\
1,512 & l.~8 & \textit{for} convergnece \textit{read} convergence \\
1,512 & l.~17 & \textit{for} clssical \textit{read} classical \\
1,512 & l.~18 & \textit{for} p.~1,343 \textit{read} p.~1,344 (Definition~6.10) \\
1,513 & l.~9 & \textit{for} p.~1,443 \textit{read} p.~1,434 (Proposition~9.7) \\
1,513 & l.~20 & \textit{for} p.~1,245 \textit{read} p.~1,244 (Lemma~3.8) \\
1,513 & l.~30 & \textit{for} vloume \textit{read} volume \\
1,514 & l.~18 & \textit{for} overveiw \textit{read} overview \\
1,514 & l.~25 & \textit{for} p.~1,467 \textit{read} p.~1,466 (Corollary~10.9) \\
1,515 & l.~18 & \textit{see} Erratum~E6 \\
1,516 & l.~18 & \textit{for} intrduction \textit{read} introduction \\
1,516 & l.~35 & \textit{for} Temkin (1993) \textit{read} Broome (1990) (the claims account of fairness) \\
1,517 & l.~6 & \textit{for} oevrview \textit{read} overview \\
1,517 & l.~7 & \textit{for} reflcetive \textit{read} reflective \\
1,517 & l.~15 & \textit{for} concepiton \textit{read} conception \\
1,517 & l.~18 & \textit{for} p.~1,438 \textit{read} p.~1,348 (Proposition~6.11) \\
1,518 & l.~15 & \textit{for} p.~1,502 \textit{read} p.~1,512 (Lemma~12.3) \\
1,518 & l.~23 & \textit{for} tohrough \textit{read} thorough \\
1,518 & l.~29 & \textit{for} guraantee \textit{read} guarantee \\
1,519 & l.~1 & \textit{for} maetrial \textit{read} material \\
1,519 & l.~4 & \textit{for} enodrsement \textit{read} endorsement \\
1,519 & l.~5 & \textit{for} sueprvenience \textit{read} supervenience \\
1,519 & l.~9 & \textit{for} persuaisve \textit{read} persuasive \\
1,519 & l.~10 & \textit{for} interprtation \textit{read} interpretation \\
1,520 & l.~29 & \textit{for} pemrissible \textit{read} permissible \\
\end{errlist}

\subsection*{Bibliography \hfill \normalfont\small pp.~1,521--1,548}
\begin{errlist}
1,522 & Anderson, E.~S. (1999) & \textit{for} A.~S. \textit{read} E.~S. \\
1,523 & Arneson, R.~J. (1989) & \textit{for} 77--92 \textit{read} 77--93 \\
1,524 & Atkinson, A.~B. (1970) & \textit{for} 234--263 \textit{read} 244--263 \\
1,525 & Broome, J. (1990) & \textit{for} 92 \textit{read} 91 \\
1,526 & Broome, J. (1991) & \textit{for} G. \textit{read} J. \\
1,526 & Cohen, G.~A. (1989) & \textit{for} 1988 \textit{read} 1989 \\
1,527 & Crisp, R. (2003) & \textit{for} 2002 \textit{read} 2003 \\
1,528 & Diamond, P.~A. (1967) & \textit{see} Erratum~E3 \\
1,528 & Dworkin, R. (1981) & \textit{for} 11(4) \textit{read} 10(4) \\
1,530 & Feinberg, J. (1970) & \textit{for} D. \textit{read} J. \\
1,531 & Frankfurt, H. (1987) & \textit{for} Equlaity \textit{read} Equality \\
1,531 & Harsanyi, J.~C. (1955) & \textit{for} comparsions \textit{read} comparisons \\
1,533 & Hirose, I. (2015) & \textit{for} F. \textit{read} I. \\
1,533 & Holtug, N. (2010) & \textit{for} 2009 \textit{read} 2010 \\
1,534 & Hooker, B. (2005) & \textit{for} A. \textit{read} B. \\
1,535 & Kagan, S. (2012) & \textit{for} Princeton University Press \textit{read} Oxford University Press \\
1,535 & Kamm, F.~M. (1993) & \textit{for} Mortlity \textit{read} Mortality \\
1,537 & Moore, G.~E. (1903) & \textit{for} H.~E. \textit{read} G.~E. \\
1,539 & Otsuka, M. and Voorhoeve, A. (2009) & \textit{for} 73(2) \textit{read} 37(2) \\
1,541 & Rawls, J. (1958) & \textit{for} B. \textit{read} J. \\
1,542 & Rawls, J. (1971) & \textit{for} L. \textit{read} J. \\
1,543 & Scanlon, T.~M. (1998) & \textit{for} Oxford University Press \textit{read} Harvard University Press \\
1,544 & Sen, A. (1973) & \textit{for} P. \textit{read} A. \\
1,544 & Sher, G. (1987) & \textit{for} Dessert \textit{read} Desert \\
1,545 & Sidgwick, H. (1907) & \textit{for} R. \textit{read} H. \\
1,546 & Stone, P. (2011) & \textit{for} Cambridge University Press \textit{read} Oxford University Press \\
1,546 & Taurek, J.~M. (1977) & \textit{for} N.~M. \textit{read} J.~M. \\
1,547 & Temkin, L.~S. (1993) & \textit{for} 1992 \textit{read} 1993 \\
1,548 & Wasserman, D. (1996) & \textit{for} chanecs \textit{read} chances \\
\end{errlist}

\subsection*{Index of names \hfill \normalfont\small pp.~1,549--1,552}
\begin{errlist}
1,550 & s.v.\ Anderson, E.~S. & \textit{for} 1,216 \textit{read} 1,226 \\
1,550 & s.v.\ Arneson, R.~J. & \textit{for} 1,405 \textit{read} 1,450 \\
1,550 & s.v.\ Atkinson, A.~B. & \textit{for} 1,245 \textit{read} 1,246 \\
1,550 & s.v.\ Barry, B. & \textit{add} 1,386 \\
1,550 & s.v.\ Broome, J. & \textit{add} 1,278 \\
1,550 & s.v.\ Crisp, R. & \textit{for} 1,434 \textit{read} 1,344 \\
1,550 & s.v.\ Diamond, P.~A. & \textit{for} 1,152 \textit{read} 1,252 \\
1,550 & s.v.\ Dworkin, R. & \textit{add} 1,245 \\
1,550 & s.v.\ Elster, J. & \textit{for} 1,248 \textit{read} 1,249 \\
1,550 & s.v.\ Feinberg, J. & \textit{for} 1,434 \textit{read} 1,435 \\
1,550 & s.v.\ Frankfurt, H. & \textit{add} 1,241 \\
1,551 & s.v.\ Harsanyi, J.~C. & \textit{add} 1,275 \\
1,551 & s.v.\ Hart, H.~L.~A. & \textit{add} 1,366 \\
1,551 & s.v.\ Hooker, B. & \textit{add} 1,169 \\
1,551 & s.v.\ McKerlie, D. & \textit{for} 1,385 \textit{read} 1,358 \\
1,551 & s.v.\ Miller, D. & \textit{for} 1,430 \textit{read} 1,431 \\
1,551 & s.v.\ Moore, G.~E. & \textit{for} 1,077 \textit{read} 1,177 \\
1,551 & s.v.\ Nozick, R. & \textit{for} 1,272 \textit{read} 1,372 \\
1,552 & s.v.\ Sen, A. & \textit{for} 1,324 \textit{read} 1,234 \\
1,552 & s.v.\ Sher, G. & \textit{for} 1,525 \textit{read} 1,425 \\
1,552 & s.v.\ Stone, P. & \textit{add} 1,258 \\
1,552 & s.v.\ Taurek, J.~M. & \textit{add} 1,489 \\
1,552 & s.v.\ Temkin, L.~S. & \textit{for} 1,245 \textit{read} 1,246 \\
1,552 & s.v.\ Wasserman, D. & \textit{for} 1,365 \textit{read} 1,265 \\
\end{errlist}

\subsection*{Index of subjects \hfill \normalfont\small pp.~1,553--1,559}
\begin{errlist}
1,554 & s.v.\ aggregation & \textit{for} 1,510 \textit{read} 1,500 \\
1,554 & s.v.\ claims & \textit{for} 1,297 \textit{read} 1,279 \\
1,554 & s.v.\ comparative desert & \textit{for} 1,455 \textit{read} 1,445 \\
1,554 & s.v.\ desert & \textit{add} 1,427 \\
1,554 & s.v.\ equality & \textit{for} 1,218 \textit{read} 1,228 \\
1,555 & s.v.\ fair play & \textit{for} 1,485 \textit{read} 1,385 \\
1,555 & s.v.\ fairness & \textit{for} 1,299 \textit{read} 1,289 \\
1,555 & s.v.\ future generations & \textit{for} 1,416 \textit{read} 1,415 \\
1,555 & s.v.\ goods of outcomes & \textit{for} 1,551 \textit{read} 1,515 \\
1,556 & s.v.\ idealized reflective endorsement & \textit{for} 1,193 \textit{read} 1,194 \\
1,556 & s.v.\ indivisible goods & \textit{for} 1,207 \textit{read} 1,217 \\
1,556 & s.v.\ isolation test & \textit{for} 1,151 \textit{read} 1,511 \\
1,556 & s.v.\ just savings & \textit{for} 1,420 \textit{read} 1,410 \\
1,557 & s.v.\ levelling-down objection & \textit{add} 1,468 \\
1,557 & s.v.\ lotteries & \textit{for} 1,089 \textit{read} 1,189 \\
1,557 & s.v.\ measurement of inequality & \textit{for} 1,237 \textit{read} 1,247 \\
1,557 & s.v.\ non-identity problem & \textit{for} 1,392 \textit{read} 1,393 \\
1,558 & s.v.\ priority view & \textit{for} 1,465 \textit{read} 1,475 \\
1,558 & s.v.\ procedural fairness & \textit{for} 1,379 \textit{read} 1,378 \\
1,558 & s.v.\ proportional satisfaction & \textit{for} 1,196 \textit{read} 1,186 \\
1,558 & s.v.\ sufficiency & \textit{for} 1,445 \textit{read} 1,455 \\
1,559 & s.v.\ supervenience & \textit{for} 1,206 \textit{read} 1,205 \\
1,559 & s.v.\ the numbers & \textit{for} 1,488 \textit{read} 1,487 \\
1,559 & s.v.\ tie-breaking argument & \textit{for} 1,399 \textit{read} 1,499 \\
1,559 & s.v.\ weighted lotteries & \textit{for} 1,375 \textit{read} 1,374 \\
\end{errlist}


\begin{thebibliography}{9}\small
\bibitem[Broome(1990)]{Broome1990} Broome, J. (1990). Fairness. \emph{Proceedings of the Aristotelian Society}, 91, 87--101.
\bibitem[Diamond(1967)]{Diamond1967} Diamond, P.~A. (1967). Cardinal welfare, individualistic ethics, and interpersonal comparison of utility: Comment. \emph{Journal of Political Economy}, 75(5), 765--766.
\bibitem[Harsanyi(1955)]{Harsanyi1955} Harsanyi, J.~C. (1955). Cardinal welfare, individualistic ethics, and interpersonal comparisons of utility. \emph{Journal of Political Economy}, 63(4), 309--321.
\bibitem[Parfit(1997)]{Parfit1997} Parfit, D. (1997). Equality and priority. \emph{Ratio}, 10(3), 202--221.
\bibitem[Temkin(1993)]{Temkin1993} Temkin, L.~S. (1993). \emph{Inequality}. Oxford University Press.
\end{thebibliography}

\end{document}
